% Options for packages loaded elsewhere
\PassOptionsToPackage{unicode}{hyperref}
\PassOptionsToPackage{hyphens}{url}
\PassOptionsToPackage{dvipsnames,svgnames,x11names}{xcolor}
\documentclass[
  11pt,
]{article}
\usepackage{xcolor}
\usepackage[margin=1in]{geometry}
\usepackage{amsmath,amssymb}
\setcounter{secnumdepth}{-\maxdimen} % remove section numbering
\usepackage{iftex}
\ifPDFTeX
  \usepackage[T1]{fontenc}
  \usepackage[utf8]{inputenc}
  \usepackage{textcomp} % provide euro and other symbols
\else % if luatex or xetex
  \usepackage{unicode-math} % this also loads fontspec
  \defaultfontfeatures{Scale=MatchLowercase}
  \defaultfontfeatures[\rmfamily]{Ligatures=TeX,Scale=1}
\fi
\usepackage[]{newtxtext}
\ifPDFTeX\else
  % xetex/luatex font selection
\fi
% Use upquote if available, for straight quotes in verbatim environments
\IfFileExists{upquote.sty}{\usepackage{upquote}}{}
\IfFileExists{microtype.sty}{% use microtype if available
  \usepackage[]{microtype}
  \UseMicrotypeSet[protrusion]{basicmath} % disable protrusion for tt fonts
}{}
\usepackage{setspace}
\makeatletter
\@ifundefined{KOMAClassName}{% if non-KOMA class
  \IfFileExists{parskip.sty}{%
    \usepackage{parskip}
  }{% else
    \setlength{\parindent}{0pt}
    \setlength{\parskip}{6pt plus 2pt minus 1pt}}
}{% if KOMA class
  \KOMAoptions{parskip=half}}
\makeatother
\usepackage{color}
\usepackage{fancyvrb}
\newcommand{\VerbBar}{|}
\newcommand{\VERB}{\Verb[commandchars=\\\{\}]}
\DefineVerbatimEnvironment{Highlighting}{Verbatim}{commandchars=\\\{\}}
% Add ',fontsize=\small' for more characters per line
\newenvironment{Shaded}{}{}
\newcommand{\AlertTok}[1]{\textcolor[rgb]{1.00,0.00,0.00}{\textbf{#1}}}
\newcommand{\AnnotationTok}[1]{\textcolor[rgb]{0.38,0.63,0.69}{\textbf{\textit{#1}}}}
\newcommand{\AttributeTok}[1]{\textcolor[rgb]{0.49,0.56,0.16}{#1}}
\newcommand{\BaseNTok}[1]{\textcolor[rgb]{0.25,0.63,0.44}{#1}}
\newcommand{\BuiltInTok}[1]{\textcolor[rgb]{0.00,0.50,0.00}{#1}}
\newcommand{\CharTok}[1]{\textcolor[rgb]{0.25,0.44,0.63}{#1}}
\newcommand{\CommentTok}[1]{\textcolor[rgb]{0.38,0.63,0.69}{\textit{#1}}}
\newcommand{\CommentVarTok}[1]{\textcolor[rgb]{0.38,0.63,0.69}{\textbf{\textit{#1}}}}
\newcommand{\ConstantTok}[1]{\textcolor[rgb]{0.53,0.00,0.00}{#1}}
\newcommand{\ControlFlowTok}[1]{\textcolor[rgb]{0.00,0.44,0.13}{\textbf{#1}}}
\newcommand{\DataTypeTok}[1]{\textcolor[rgb]{0.56,0.13,0.00}{#1}}
\newcommand{\DecValTok}[1]{\textcolor[rgb]{0.25,0.63,0.44}{#1}}
\newcommand{\DocumentationTok}[1]{\textcolor[rgb]{0.73,0.13,0.13}{\textit{#1}}}
\newcommand{\ErrorTok}[1]{\textcolor[rgb]{1.00,0.00,0.00}{\textbf{#1}}}
\newcommand{\ExtensionTok}[1]{#1}
\newcommand{\FloatTok}[1]{\textcolor[rgb]{0.25,0.63,0.44}{#1}}
\newcommand{\FunctionTok}[1]{\textcolor[rgb]{0.02,0.16,0.49}{#1}}
\newcommand{\ImportTok}[1]{\textcolor[rgb]{0.00,0.50,0.00}{\textbf{#1}}}
\newcommand{\InformationTok}[1]{\textcolor[rgb]{0.38,0.63,0.69}{\textbf{\textit{#1}}}}
\newcommand{\KeywordTok}[1]{\textcolor[rgb]{0.00,0.44,0.13}{\textbf{#1}}}
\newcommand{\NormalTok}[1]{#1}
\newcommand{\OperatorTok}[1]{\textcolor[rgb]{0.40,0.40,0.40}{#1}}
\newcommand{\OtherTok}[1]{\textcolor[rgb]{0.00,0.44,0.13}{#1}}
\newcommand{\PreprocessorTok}[1]{\textcolor[rgb]{0.74,0.48,0.00}{#1}}
\newcommand{\RegionMarkerTok}[1]{#1}
\newcommand{\SpecialCharTok}[1]{\textcolor[rgb]{0.25,0.44,0.63}{#1}}
\newcommand{\SpecialStringTok}[1]{\textcolor[rgb]{0.73,0.40,0.53}{#1}}
\newcommand{\StringTok}[1]{\textcolor[rgb]{0.25,0.44,0.63}{#1}}
\newcommand{\VariableTok}[1]{\textcolor[rgb]{0.10,0.09,0.49}{#1}}
\newcommand{\VerbatimStringTok}[1]{\textcolor[rgb]{0.25,0.44,0.63}{#1}}
\newcommand{\WarningTok}[1]{\textcolor[rgb]{0.38,0.63,0.69}{\textbf{\textit{#1}}}}
\usepackage{longtable,booktabs,array}
\newcounter{none} % for unnumbered tables
\usepackage{calc} % for calculating minipage widths
% Correct order of tables after \paragraph or \subparagraph
\usepackage{etoolbox}
\makeatletter
\patchcmd\longtable{\par}{\if@noskipsec\mbox{}\fi\par}{}{}
\makeatother
% Allow footnotes in longtable head/foot
\IfFileExists{footnotehyper.sty}{\usepackage{footnotehyper}}{\usepackage{footnote}}
\makesavenoteenv{longtable}
% definitions for citeproc citations
\NewDocumentCommand\citeproctext{}{}
\NewDocumentCommand\citeproc{mm}{%
  \begingroup\def\citeproctext{#2}\cite{#1}\endgroup}
\makeatletter
 % allow citations to break across lines
 \let\@cite@ofmt\@firstofone
 % avoid brackets around text for \cite:
 \def\@biblabel#1{}
 \def\@cite#1#2{{#1\if@tempswa , #2\fi}}
\makeatother
\newlength{\cslhangindent}
\setlength{\cslhangindent}{1.5em}
\newlength{\csllabelwidth}
\setlength{\csllabelwidth}{3em}
\newenvironment{CSLReferences}[2] % #1 hanging-indent, #2 entry-spacing
 {\begin{list}{}{%
  \setlength{\itemindent}{0pt}
  \setlength{\leftmargin}{0pt}
  \setlength{\parsep}{0pt}
  % turn on hanging indent if param 1 is 1
  \ifodd #1
   \setlength{\leftmargin}{\cslhangindent}
   \setlength{\itemindent}{-1\cslhangindent}
  \fi
  % set entry spacing
  \setlength{\itemsep}{#2\baselineskip}}}
 {\end{list}}
\usepackage{calc}
\newcommand{\CSLBlock}[1]{\hfill\break\parbox[t]{\linewidth}{\strut\ignorespaces#1\strut}}
\newcommand{\CSLLeftMargin}[1]{\parbox[t]{\csllabelwidth}{\strut#1\strut}}
\newcommand{\CSLRightInline}[1]{\parbox[t]{\linewidth - \csllabelwidth}{\strut#1\strut}}
\newcommand{\CSLIndent}[1]{\hspace{\cslhangindent}#1}
\setlength{\emergencystretch}{3em} % prevent overfull lines
\providecommand{\tightlist}{%
  \setlength{\itemsep}{0pt}\setlength{\parskip}{0pt}}
\usepackage{amsmath,amssymb,amsthm}
\usepackage{newtxmath}
\usepackage{booktabs}
\usepackage{threeparttable}
\usepackage{graphicx}
\usepackage{float}
\usepackage{pdflscape}
\usepackage{needspace}
\usepackage{setspace}
\usepackage{etoolbox}
\usepackage[font=small,labelfont=bf,labelsep=period]{caption}
\usepackage[section]{placeins}
\definecolor{paperlink}{RGB}{31,59,115}
\allowdisplaybreaks
\setlength{\parskip}{0pt}
\setlength{\parindent}{1.5em}
\AtBeginEnvironment{CSLReferences}{\interlinepenalty=10000}
\usepackage{fvextra}
\fvset{breaklines=true,breakanywhere=true}
\DefineVerbatimEnvironment{verbatim}{Verbatim}{breaklines=true,breakanywhere=true}
\AtBeginEnvironment{longtable}{\footnotesize}
\usepackage{bookmark}
\IfFileExists{xurl.sty}{\usepackage{xurl}}{} % add URL line breaks if available
\urlstyle{same}
\hypersetup{
  pdftitle={Online Appendix to Competition Creates Competition},
  pdfauthor={Austin Li},
  colorlinks=true,
  linkcolor={paperlink},
  filecolor={Maroon},
  citecolor={paperlink},
  urlcolor={paperlink},
  pdfcreator={LaTeX via pandoc}}

\ifXeTeX
\special{pdf:trailerid [ <c0fc91c2f063f65525272c64d93d2030> <c0fc91c2f063f65525272c64d93d2030> ]}
\fi
\ifPDFTeX
\pdftrailerid{}
\pdftrailer{/ID [ <c0fc91c2f063f65525272c64d93d2030> <c0fc91c2f063f65525272c64d93d2030> ]}
\fi
\ifLuaTeX
\pdfvariable trailerid {[ <c0fc91c2f063f65525272c64d93d2030> <c0fc91c2f063f65525272c64d93d2030> ]}
\fi

\title{Online Appendix to Competition Creates Competition}
\author{Austin Li}
\date{}

\begin{document}
\maketitle

\setstretch{1.15}
\section{A. Full analytical arguments}\label{oa-a}

This appendix develops the probability space, proves the main results
and extensions, and states the regularity conditions used in the
sale-design calculations. Sections B and C describe the equilibrium
certificates and numerical exercises.

\subsection{A.1. Probability space, strategies, and conditional
information}\label{oa-a-foundations}

The benchmark has a risk-neutral seller, an incumbent, a potential
challenger, a strategic investor, competitive market makers, and
exogenous noise demand. A known target standalone value is normalized to
zero. The seller commits before trading to a binding second-price cash
auction with reserve \(p\). The incumbent has value \(R\sim U[0,r]\),
pays no incremental participation cost at the modeled stage, and learns
its value before bidding. The challenger has value \(\ell\) or \(h\)
with equal probability, where \(0<p<\ell<r<h\). Neither bidder trades
target equity. An outside investor observes challenger quality and
chooses an order in \([-1,1]\), pays \(k|q|\), and earns
\(q(V_T-P)-k|q|\). Its order unit is not a controlling ownership stake.
All payoffs and preparation costs use the same per-share normalization.

Noise has density \(f(z)=e^{-|z|/b}/(2b)\), with \(b>1\) and \(k>0\).
Market makers observe \(X=q+Z\) and price \(P(X)=\mathbb E[V_T\mid X]\).
The challenger observes \(P\) but not \(X\), then its independent
preparation cost, equal to \(c_L\) with probability \(\rho\) and \(c_H\)
otherwise, where \(0<\rho<1\) and \(0\le c_L<c_H\). Paying the cost
reveals its value and permits bidding; declining means absence. Both
bidders bid truthfully once informed. A bid meeting the reserve is
admissible. I prescribe preparation at equality of gross expected profit
and cost. The allocation, cash payment, and stock payoff are then
realized. The equilibrium requires Bayesian pricing, optimal preparation
at every reached price and cost, and global investor best responses
against the fixed equilibrium price and preparation schedules,
permitting mixed orders.

I write \(t_0\) for expected target proceeds without challenger entry,
\(t_H,t_L\) for proceeds conditional on entry and quality, and
\(g_H,g_L\) for conditional gross challenger profits. Set
\(\Delta_T=t_H-t_L\), \(B_r(\mu)=g_L+\mu(g_H-g_L)\),
\(m=(1+e^{2/b})^{-1}\), and \(M=1-m\). For the comparison
\(p<\ell<r_0<r_1<h\), the hypotheses of Proposition 2, referred to below
as (A1) to (A3), are

\[
\begin{aligned}
&0\le c_L<B_{r_1}(m),\\
&B_{r_0}(1/2)<c_H<B_{r_1}(M),\\
&\Delta_T(r_0)<k<(1-1/b)\rho m\Delta_T(r_1).
\end{aligned}
\tag{OA.1}\label{eq:oa-oa-1}
\]

To define beliefs and deviations, let \(\Theta\in\{L,H\}\) have equal
probabilities, with acquisition values \(v_L=\ell\) and \(v_H=h\).
Independently draw incumbent value \(R\), preparation cost \(C\), noise
demand \(Z\), and an auxiliary random variable \(U\) uniform on
\([0,1]\). A mixed trader strategy is a probability kernel
\(\sigma_\Theta\) on the Borel subsets of \([-1,1]\). It can be
implemented by a Borel quantile map \(q_\Theta(U)\). Conditional
randomization therefore contains no information about the other
primitives beyond the trader's specified signal.

All spaces are standard Borel. Prices, entry policies, and payoffs are
Borel functions. Because values and the order interval are bounded, all
acquisition and trading payoffs are integrable. A regular conditional
distribution of quality given price exists. Expectations conditional on
a price atom use its complete preimage, not a pointwise inversion that
presumes injectivity.

For the benchmark, define

\[
\begin{aligned}
a_\theta(x)&=\int_{[-1,1]}f(x-q)\,\sigma_\theta(dq),\\
\nu(dx)&=\frac{a_H(x)+a_L(x)}2\,dx,\\
\mu_X(x)&=\frac{a_H(x)}{a_H(x)+a_L(x)}.
\end{aligned}
\tag{OA.2}\label{eq:oa-oa-2}
\]

The density \(f\) is positive everywhere, bounded, continuous, and
integrates to unity. Tonelli's theorem gives \(\int a_\theta=1\).
Dominated convergence in \(q\) implies continuity of \(a_\theta(x)\),
and positivity gives a continuous posterior. The distribution of \(X\)
under every pure or mixed unilateral order deviation is absolutely
continuous with respect to Lebesgue measure. Moreover, \(\nu\) is
equivalent to Lebesgue measure because its density is strictly positive.
Thus an equilibrium identity holding \(\nu\)-almost surely also holds
almost surely under every unilateral trading deviation. No deviation can
exploit a different version of prices or beliefs on a null set.

Write \(\mathcal G=\sigma(P(X))\) and
\(\mu_P=\mathbb E[\mathbf1\{\Theta=H\}\mid\mathcal G]\). Since \(P(X)\)
is measurable with respect to \(X\),

\[
\mu_P=\mathbb E[\mu_X(X)\mid\mathcal G].
\tag{OA.3}\label{eq:oa-oa-3}
\]

For any \(q,q'\in[-1,1]\), the triangle inequality implies

\[
\left||x-q|-|x-q'|\right|\le|q-q'|\le2.
\tag{OA.4}\label{eq:oa-oa-4}
\]

Exponentiating and integrating with respect to
\(\sigma_H(dq)\sigma_L(dq')\) gives
\(e^{-2/b}a_L\le a_H\le e^{2/b}a_L\), hence \(m\le\mu_X\le M\). Equation
\eqref{eq:oa-oa-3} preserves these bounds for \(\mu_P\). These
conclusions do not impose pure trading, monotone orders, or an
informative price.

The bidder's preparation rule is \(\mathbf1\{C\le B_r(\mu_P)\}\), with
entry at equality. When \(c_L<B_r(m)\), the cost-averaged rule is
bounded below by \(\rho\). Conditional on \(\Theta\) and \(X\), the
variables \(R,C\) remain independent of the investor's action and retain
their specified laws. This is why the auction-stage expectations can be
used inside the pricing equation. It also explains why a unilateral
order deviation shifts the density of \(X\) without changing the
conditional auction formulas.

For the complementary-signal extension, I replace the investor's
observation of \(\Theta\) by \(T\) and add \(Y\). The randomization
\(U\), noise \(Z\), incumbent \(R\), and cost \(C\) remain independent
of \((\Theta,T,Y)\); \(T\) and \(Y\) are independent conditional on
\(\Theta\). Sections A.7 and C.3 specify the resulting conditional laws
rather than treating the buyer's private signal as public.

\subsection{A.2. Auction implementation and Proposition
1}\label{oa-a-payoffs}

Conditional on the highest competing admissible bid, truthful bidding
maximizes a private-value bidder's payoff in a second-price auction. A
bid above value can turn a loss into an unprofitable win, and a bid
below value can discard a profitable win. Conditional on winning, the
bidder's own bid does not determine payment. The presence of a public
price or information conveyed by entry changes beliefs about other
bidders but not this pointwise dominance argument. I use truthful
implementation to fix payoff-equivalent bid descriptions.

Without the challenger, the seller receives \(p\) if \(R\ge p\). A
high-value challenger wins almost surely and pays \(\max(p,R)\) because
\(h\) exceeds incumbent support. When the challenger has value \(\ell\),
the winning value is \(\max(R,\ell)\) and the runner-up is
\(\min(R,\ell)\), so seller payment is \(\max(p,\min(R,\ell))\). The
challenger obtains \((\ell-\max(p,R))_+\).

The uniform formulas follow by splitting at \(p\) and \(\ell\):

\[
\begin{aligned}
t_0&=\frac1r\int_p^r p\,du,\\
t_H&=\frac1r\left[p^2+\frac{r^2-p^2}{2}\right],\\
t_L&=\frac1r\left[p^2+\frac{\ell^2-p^2}{2}+\ell(r-\ell)\right],\\
g_L&=\frac1r\left[p(\ell-p)+\frac{(\ell-p)^2}{2}\right],\qquad g_H=h-t_H.
\end{aligned}
\tag{OA.5}\label{eq:oa-oa-5}
\]

Evaluating these expressions gives \(t_0=p(1-p/r)\),
\(t_H=r/2+p^2/(2r)\), \(t_L=\ell-(\ell^2-p^2)/(2r)\), \(g_H=h-t_H\), and
\(g_L=(\ell^2-p^2)/(2r)\). Subtracting gives
\(\Delta_T=(r-\ell)^2/(2r)\). On \(p<\ell<r<h\), the derivatives are
\(\Delta_T'=1/2-\ell^2/(2r^2)>0\), \(g_H'=-1/2+p^2/(2r^2)<0\), and
\(g_L'=-(\ell^2-p^2)/(2r^2)<0\). Thus \(\partial_rB_r(\mu)<0\) at every
fixed posterior. In particular \(g_H-g_L>0\), since the high challenger
has a strictly higher value in a mechanism whose allocation is monotone
in its own value.

For a general continuous incumbent CDF \(F\) on \([0,\bar r]\), the
payment difference is zero for \(R\le\ell\) and \(R-\ell\) for
\(R>\ell\). By writing a positive part as an integral of indicators,

\[
(R-\ell)_+=\int_\ell^{\bar r}\mathbf1\{R>u\}\,du,
\qquad
(\theta-\max(p,R))_+=\int_p^\theta\mathbf1\{R<u\}\,du.
\tag{OA.6}\label{eq:oa-oa-6}
\]

Tonelli applies because the integrands are nonnegative and the domains
bounded. CDF continuity removes the distinction between \(R<u\) and
\(R\le u\) in expectation; even without continuity the Lebesgue integral
over \(u\) is unaffected by countably many atoms. Thus the two
expressions are respectively \(\int_\ell^{\bar r}(1-F(u))du\) and
\(\int_p^\theta F(u)du\), extending \(F=1\) above support. A first-order
stochastic strengthening lowers \(F\) pointwise. Nonnegative integral
differences prove weak ordering; a strictly positive integral difference
proves strictness. Posterior-weighted challenger profits inherit the
ordering. This completes Proposition 1.

\subsection{A.3. Price sufficiency, residuals, and convolution
regularity}\label{oa-a-inference}

For any candidate equilibrium with the low-cost entry floor, the
independent cost distribution gives a Borel function \(e(P)\ge\rho\).
Conditional on \(X=x\), expected target payoff is

\[
P(x)=t_0+e(P(x))\{t_L-t_0+\Delta_T\mu_X(x)\}.
\tag{OA.7}\label{eq:oa-oa-7}
\]

The function \(g(P)=[P-t_0-e(P)(t_L-t_0)]/[e(P)\Delta_T]\) is Borel on
the realized price range because its denominator is strictly positive.
It satisfies \(\mu_X(X)=g(P(X))\) almost surely. Equation
\eqref{eq:oa-oa-3} then gives \(\mu_P=\mu_X\) almost surely. The entry
function is consequently

\[
e(\mu)=\rho+(1-\rho)\mathbf1\{g_L+\mu(g_H-g_L)\ge c_H\}.
\tag{OA.8}\label{eq:oa-oa-8}
\]

For construction, put \(w_L=t_L-t_0>0\). For \(\mu_2>\mu_1\),

\[
\begin{aligned}
P(\mu_2)-P(\mu_1)
&=e(\mu_2)\Delta_T(\mu_2-\mu_1)
+[e(\mu_2)-e(\mu_1)](w_L+\Delta_T\mu_1)\\
&\ge\rho\Delta_T(\mu_2-\mu_1)>0.
\end{aligned}
\tag{OA.9}\label{eq:oa-oa-9}
\]

Thus \(P(\mu)\) is strictly increasing even with an entry jump. Its
inverse on its image is measurable: it is the restriction of a
generalized inverse defined through infima of upper level sets. Price
values inside a skipped interval are not generated by any flow and need
not receive artificial posterior mass. By contrast, a flat posterior
tail maps into an actual price atom and is conditioned on as an atom.

Conditional on fundamental quality, expected terminal value under a
unilateral trading deviation is still \(t_0+e(P(x))(t_\theta-t_0)\).
Subtracting \eqref{eq:oa-oa-7} gives \(A_H=e\Delta_T(1-\mu_X)\) and
\(A_L=e\Delta_T\mu_X\). Both residuals are bounded in
\([\rho m\Delta_T,\Delta_T]\). The signs exclude incorrectly signed
orders. This derivation accounts for the endogenous preparation response
before taking the investor's expectation over shifted noise.

Next I justify every differentiation used for continuous deviations. Let
\(f\in W^{1,1}(\mathbb R)\), with positive density and \(|f'|\le Lf\)
almost everywhere, and let \(A\) be bounded, measurable, and
nonnegative. For sign \(\epsilon\in\{-1,1\}\), define

\[
F_\epsilon(s)=\int_{\mathbb R}f(x-\epsilon s)A(x)\,dx.
\tag{OA.10}\label{eq:oa-oa-10}
\]

For \(s_1<s_2\), the absolutely continuous representative of \(f\)
satisfies

\[
f(x-\epsilon s_2)-f(x-\epsilon s_1)
=-\epsilon\int_{s_1}^{s_2}f'(x-\epsilon u)\,du.
\tag{OA.11}\label{eq:oa-oa-11}
\]

The absolute integral of the right-hand side after multiplication by
\(A\) is at most \(\|A\|_\infty(s_2-s_1)\|f'\|_1\). Fubini therefore
permits interchanging the integrals. This proves absolute continuity of
\(F_\epsilon\) and

\[
F_\epsilon'(s)=-\epsilon\int f'(x-\epsilon s)A(x)\,dx,
\qquad |F_\epsilon'(s)|\le L F_\epsilon(s)
\tag{OA.12}\label{eq:oa-oa-12}
\]

almost everywhere. Translation continuity of \(f'\) in \(L^1\) also
makes the displayed derivative continuous. One way to see the needed
translation continuity without an additional smoothness assumption is to
approximate \(f'\) in \(L^1\) by compactly supported continuous
functions, for which it holds uniformly; the approximation error is
unchanged by translation. The difference quotient of \eqref{eq:oa-oa-11}
then converges in \(L^1\) to \(-\epsilon f'(\cdot-\epsilon s)\),
identifying the classical derivative everywhere. The almost-everywhere
statement is already sufficient for the proof.

The Laplace density is in \(W^{1,1}\): it is continuous at zero,
absolutely continuous on every bounded interval, and has integrable weak
derivative \(-\operatorname{sgn}(x)f(x)/b\). Its kink contributes no
atom to the first weak derivative because the density itself is
continuous. Hence \(L=1/b\). This argument differentiates the translated
density, not the discontinuous entry policy.

For \(U(s)=sF(s)-ks\), absolute continuity and the product rule imply

\[
U'(s)=F(s)+sF'(s)-k\ge(1-sL)F(s)-k.
\tag{OA.13}\label{eq:oa-oa-13}
\]

If the right-hand side is bounded below by a strictly positive constant
on the allowed interval, integration shows that \(U(s_2)>U(s_1)\) for
every \(s_2>s_1\). This is global order optimality, not a local
first-order condition.

\subsection{A.4. Complete proof of Proposition 2}\label{oa-a-entry}

The proof first restricts the trading strategy in any candidate
equilibrium, then constructs prices and entry that support it.

\textbf{Weak incumbent.} Proposition A.1 and (A1) give the entry floor
under every conditional mixed order. \hyperref[oa-a-inference]{Section
A.3} gives the residual upper bound. For a correctly signed order of
magnitude \(s>0\),

\[
\mathbb E[\text{trading profit}\mid\Theta]\le s\{\Delta_T(r_0)-k\}<0.
\tag{OA.14}\label{eq:oa-oa-14}
\]

An incorrectly signed order has strictly negative gross expected profit.
Zero gives zero. Thus each type has the unique best response zero,
irrespective of the candidate equilibrium schedules. A mixed
distribution placing positive probability on a nonzero order cannot be
optimal: its payoff is an average of strictly negative payoffs and zero,
and is strictly negative if that probability is positive. There is no
issue from arbitrarily small trades; a nonpositive integrable random
payoff that is strictly negative on a positive-probability set has a
negative expectation.

Under zero orders, \(X=Z\) is independent of quality. Set \(e=\rho\) and
\(P=t_0+\rho[(t_H+t_L)/2-t_0]\). Condition (A2) excludes high-cost
preparation at the prior, and (A1) admits low-cost preparation.
Competitive pricing is correct and the previous payoff calculation
validates zero. This constructs the unique trading and on-path
preparation outcome.

\textbf{Strong incumbent.} Again consider any candidate equilibrium,
including mixed orders. \hyperref[oa-a-inference]{Section A.3} gives
\(F_\theta(s)\ge\rho m\Delta_T(r_1)\), and \eqref{eq:oa-oa-13} with
\(L=1/b\) gives

\[
U_\theta'(s)\ge(1-1/b)\rho m\Delta_T(r_1)-k>0.
\tag{OA.15}\label{eq:oa-oa-15}
\]

Every smaller correctly signed magnitude is strictly inferior to the
unit magnitude; incorrect signs are inferior to zero. Thus the only
possible equilibrium orders are \(q_H=1,q_L=-1\).

Their conditional flow densities are \(f(x-1)\) and \(f(x+1)\). Bayes'
rule gives the piecewise posterior

\[
\mu(x)=
\begin{cases}
m,&x\le-1,\\
(1+e^{-2x/b})^{-1},&-1<x<1,\\
M,&x\ge1.
\end{cases}
\tag{OA.16}\label{eq:oa-oa-16}
\]

Use the preparation rule \eqref{eq:oa-oa-8} and pricing
\eqref{eq:oa-oa-7}. Equation \eqref{eq:oa-oa-9} makes the price
invertible with respect to this posterior, so the buyer optimizes using
exactly the information it observes. The global bounds verify investor
best responses, and the construction satisfies competitive pricing and
truthful bidding. This proves existence and the uniqueness of the
trading and on-path preparation outcome. Price functions can differ at
unrealized price values or flow null sets without creating a different
economic outcome.

\textbf{Entry and allocation.} Since \(g_H>g_L\), expensive preparation
requires \(\mu\ge\tau\). The strict inequalities
\(B_{r_1}(1/2)<B_{r_0}(1/2)<c_H<B_{r_1}(M)\) imply \(\tau\in(1/2,M)\).
Solving \eqref{eq:oa-oa-16} gives
\(x^*=b\operatorname{logit}(\tau)/2\in(0,1)\). The Laplace survival
function is

\[
\overline F_Z(z)=
\begin{cases}
1-\tfrac12e^{z/b},&z<0,\\
\tfrac12e^{-z/b},&z\ge0.
\end{cases}
\tag{OA.17}\label{eq:oa-oa-17}
\]

Hence \(\alpha_H=\overline F_Z(x^*-1)\) and
\(\alpha_L=\overline F_Z(x^*+1)\) give the main-text formulas. Both are
positive. State-specific entry is
\(e_\theta=\rho+(1-\rho)\alpha_\theta\), total entry is \((e_H+e_L)/2\),
and high-quality ownership is \(e_H/2\) because \(h>r\). Each comparison
is strict.

\textbf{Information ordering.} A signal experiment is its pair of
conditional distributions given quality. Applying a constant Markov
kernel to the strong-economy price produces the weak-economy constant
price experiment. Conversely, every state-independent kernel applied to
a constant experiment has the same conditional law in both states. The
strong experiment does not: on an interior price region its posterior is
nonconstant, since the conditional flow likelihood ratio is nonconstant
and the price reveals it. Thus no reverse garbling exists. Differences
in mean payoffs across economies do not affect this comparison of signal
experiments.

\textbf{Nonemptiness.} \hyperref[oa-a-extensions]{Section A.6} gives a
direct primitive construction for every \(h>\ell\). The specified
benchmark is also strictly inside the condition region. All inequalities
involve continuous primitive expressions away from the support
boundaries, so their strict satisfaction persists on an open
neighborhood.

\subsection{A.5. Fixed experiments, threshold distinctions, and finite
nonmonotonicity}\label{oa-a-thresholds}

For Proposition A.3, fix one joint distribution of \((S,\Theta)\) for
both strengths, and keep the preparation cost \(C\) independent of
\((S,\Theta)\) with its law also fixed as strength varies. Only \(r\)
changes between the two economies; the signal experiment and the cost
law do not. Cost independence is what makes the buyer's posterior
\(\mu(S)=\Pr(H\mid S)\) the right object: with a cost that carried
information about quality, the comparison would need \(\Pr(H\mid S,C)\)
and an independence condition between the incumbent's value and the
buyer's full information, and I do not claim that version. For
\(r_1>r_0\), the difference in entry indicators is

\[
\mathbf1\{C\le B_{r_0}(\mu(S))\}
-\mathbf1\{C\le B_{r_1}(\mu(S))\}
=\mathbf1\{B_{r_1}(\mu(S))<C\le B_{r_0}(\mu(S))\}.
\tag{OA.18}\label{eq:oa-oa-18}
\]

Taking expectations proves weak decrease and the exact strictness
condition, without requiring an atomless cost distribution. This
comparison holds within a fixed full-order regime because both
conditional flow laws are unchanged when only \(r\) changes. It need not
hold across different informative order profiles.

For Proposition A.4, define the maintained domain

\[
\mathcal D=\{r\in(\ell,h):c_L<B_r(m),\ B_r(1/2)<c_H\}.
\tag{OA.19}\label{eq:oa-oa-19}
\]

Under pooling, the public posterior is \(1/2\) and entry is \(\rho\). A
unilateral correctly signed order sees a constant residual
\(\rho\Delta_T/2\) and earns \(s(\rho\Delta_T/2-k)\). Thus pooling is an
equilibrium exactly when this coefficient is nonpositive. Equality
supports pooling although the trader is indifferent among correctly
signed sizes against that constant candidate schedule; it does not make
every alternative profile an equilibrium with recomputed prices.

Solving \(\Delta_T(r)=d\) gives \((r-\ell)^2=2rd\), whose roots are
\(\ell+d\pm\sqrt{d^2+2\ell d}\). The plus root is the one above
\(\ell\). Writing \(\mathfrak r(d)=\ell+d+\sqrt{d^2+2\ell d}\) gives the
exact pooling-existence boundary \(r_N=\mathfrak r(2k/\rho)\), the
sufficient pooling-uniqueness boundary \(\mathfrak r(k)\), and the
sufficient full-order uniqueness boundary
\(r_U=\mathfrak r(k/[(1-1/b)\rho m])\). These boundaries are applied
only within \(\mathcal D\). A sufficient global bound does not identify
the first existence of an informative equilibrium.

At \(r>r_U\) in \(\mathcal D\), the global derivative bound forces full
orders. Expensive entry is feasible only if the attainable posterior
reaches the required \(\tau\). The uniform mixed-strategy posterior
bound makes \(B_r(M)<c_H\) sufficient to exclude expensive entry in
every equilibrium, regardless of the orders. The relevant equality is

\[
Mr^2-2(Mh-c_H)r-[(1-M)\ell^2-p^2]=0.
\tag{OA.20}\label{eq:oa-oa-20}
\]

Because \(B_r(M)\) is strictly decreasing on \(r>\ell>p\), there is at
most one root in the economic domain. The larger quadratic root gives
that root when it lies in the domain; if both algebraic roots are
positive, the other cannot lie in this strictly decreasing domain as a
second crossing. Numerical specifications must verify the discriminant,
the root residual, and the support and floor conditions rather than
treating the formula as globally applicable.

For Laplace noise, the posterior \(M\) is attained on \(x\ge1\). At
equality \(B_r(M)=c_H\), entry-at-indifference admits the costly buyer
on that complete tail. As \(\tau\uparrow M\), the flow threshold
approaches unity and

\[
\alpha_H\to\frac12,\qquad \alpha_L\to\frac12e^{-2/b},\qquad
\mathsf E\to\rho+\frac{1-\rho}{4}(1+e^{-2/b}).
\tag{OA.21}\label{eq:oa-oa-21}
\]

For any strict infeasibility beyond the crossing, costly entry is zero.
An alternative tie rule changes the value at the equality itself, not
the strict comparisons on either side. With logistic noise the upper
posterior is approached only as \(x\to\infty\); its survival probability
tends to zero, and entry converges continuously to \(\rho\).

For Proposition 2(iii), impose (A1) to (A3) at \(r_0,r_1\), the low-cost
floor at \(r_2\), the strong global derivative bound at \(r_2\), and
\(B_{r_2}(M)<c_H\). Parts (i) and (ii) of Proposition 2 give the first
two unique outcomes; the bound and infeasibility condition give unique
full orders but entry \(\rho\) at \(r_2\). This proves a finite rise and
fall without a selection or uniqueness assertion for strengths between
them.

A useful full-profile existence test is sharper than the uniform
uniqueness bound. At a fixed full-order candidate,

\[
J=F_H(1)=F_L(1)
=\Delta_T\int e(x)\frac{f(x-1)f(x+1)}{f(x-1)+f(x+1)}\,dx.
\tag{OA.22}\label{eq:oa-oa-22}
\]

The equality follows by substituting Bayes' rule into both residuals.
The pointwise density-ratio bound implies
\(F_\theta(s)\ge e^{-(1-s)/b}J\). Since \((1-s/b)e^{-(1-s)/b}\) is
decreasing on \([0,1]\), \(k<(1-1/b)J\) makes every correctly signed
marginal profit positive. It verifies existence of that profile, not
uniqueness over other candidate schedules. The distinction is preserved
in the numerical classifications.

\subsection{A.6. Smooth noise, atomless costs, and primitive
nonemptiness}\label{oa-a-extensions}

For logistic noise,

\[
f(z)=\frac{e^{-z/b}}{b(1+e^{-z/b})^2},\qquad
(\log f)'(z)=-\frac1b\tanh\left(\frac z{2b}\right).
\tag{OA.23}\label{eq:oa-oa-23}
\]

It is positive and smooth, and \(|f'|\le f/b\) implies
\(\|f'\|_1\le1/b\). Therefore it belongs to \(W^{1,1}\) and all
posterior and global trading bounds apply. Under full orders the
posterior log odds equal

\[
\log\frac{\mu(x)}{1-\mu(x)}
=2\log\cosh\left(\frac{x+1}{2b}\right)
-2\log\cosh\left(\frac{x-1}{2b}\right).
\tag{OA.24}\label{eq:oa-oa-24}
\]

Their derivative equals

\[
\frac1b\left[\tanh\left(\frac{x+1}{2b}\right)-\tanh\left(\frac{x-1}{2b}\right)\right]>0.
\tag{OA.25}\label{eq:oa-oa-25}
\]

The limits are \(-2/b\) and \(2/b\), giving the open posterior range
\((m,M)\). Every strictly interior threshold is crossed with positive
probability. To derive its inverse, set \(y=e^{x/b}\), \(A=e^{1/b}\),
and \(w=\sqrt{\tau/(1-\tau)}\). The square root of the likelihood ratio
equals \((Ay+1)/(y+A)\). Solving \(w=(Ay+1)/(y+A)\) gives
\(y=(Aw-1)/(A-w)>0\). Thus \(x^*_{\log}=b\log[(Aw-1)/(A-w)]\). Its
conditional survival probabilities are
\(\alpha_H=[1+e^{(x^*_{\log}-1)/b}]^{-1}\) and
\(\alpha_L=[1+e^{(x^*_{\log}+1)/b}]^{-1}\). These formulas and
\(\mathsf E=\rho+(1-\rho)(\alpha_H+\alpha_L)/2\) supply the complete
entry comparison. The trade and information-ordering arguments complete
Proposition A.5.

For Proposition A.6, let \(H_C=\rho H_L+(1-\rho)H_H\), with atomless
component laws supported within
\([c_L-\varepsilon_C,c_L+\varepsilon_C]\) and
\([c_H-\varepsilon_C,c_H+\varepsilon_C]\). Retain (A3) and impose
\(c_L-\varepsilon_C\ge0\), \(c_L+\varepsilon_C<B_{r_1}(m)\), and

\[
B_{r_0}(1/2)<c_H-\varepsilon_C<c_H+\varepsilon_C<B_{r_1}(M).
\tag{OA.26}\label{eq:oa-oa-26}
\]

Use either the Laplace or logistic noise law. The entry rule is
\(e(\mu)=H_C(B_r(\mu))\). It is continuous and nondecreasing, and
\(e\ge\rho\) under the low-cost support restriction. Equation
\eqref{eq:oa-oa-9} remains strictly positive even if the component
densities are discontinuous or vanish on subintervals. In the weak
economy, \(e(1/2)=\rho\). In the strong economy, a posterior
sufficiently close to \(M\) satisfies \(B_r(\mu)>c_H+\varepsilon_C\); an
event with positive probability then gives \(e=1\). All global trading
comparisons are unchanged. This proof requires an atomless distribution,
not a differentiable density. Differentiability is imposed separately
when differentiating a seller objective.

For primitive nonemptiness, fix \(h>\ell>p>0\), \(b>1\), and
\(0<\rho<1\). At the continuous limit \(r=\ell\),

\[
g_H-g_L=h-\ell>0,\qquad
B_\ell(M)-B_\ell(1/2)=(M-1/2)(h-\ell)>0.
\tag{OA.27}\label{eq:oa-oa-27}
\]

Choose \(r_1\in(\ell,h)\) close enough to \(\ell\) that
\(B_{r_1}(M)>B_\ell(1/2)\). Because \(B_r(1/2)\) decreases in \(r\),
every \(r_0\in(\ell,r_1)\) then satisfies \(B_{r_0}(1/2)<B_{r_1}(M)\).
Choose \(r_0\) sufficiently near \(\ell\) that
\(\Delta_T(r_0)<(1-1/b)\rho m\Delta_T(r_1)\). Select \(k\) strictly
inside this interval and \(c_H\) strictly between the two profit bounds.
The quantity \(B_{r_1}(m)\) is positive, so choose \(c_L\) strictly
between zero and that bound. The constructed \(c_L<c_H\) follows because
\(B_{r_1}(m)<B_{r_0}(1/2)\), using \(m<1/2\) and the monotonicity in
both arguments. All inequalities have strict slack and persist by
continuity. This is a mathematical nonemptiness construction, not a
statement about empirical effect size at every value ratio.

Two further stability results are recorded for completeness.

\Needspace{8\baselineskip}

\protect\phantomsection\label{result-lemma-oa-1}{\textbf{Lemma OA.1
(analytical).}} \emph{Replace the trading cost by
\(K(s)=ks+\gamma_qs^2/2\), \(\gamma_q\ge0\). Under (A1) and (A2),
\(\Delta_T(r_0)<k\), and \(k+\gamma_q<(1-1/b)\rho m\Delta_T(r_1)\), the
unique trading outcomes and entry reversal in parts (i) and (ii) of
Proposition 2 remain valid.}

Indeed \(K(s)\ge ks\) preserves the weak-economy exclusion. At high
strength, subtracting \(K'(s)=k+\gamma_qs\) from the gross marginal
profit leaves a strictly positive bound. The price construction,
threshold, and information comparison are unchanged. This perturbation
retains the position bound; it does not solve an unbounded-order game.

\Needspace{8\baselineskip}

\protect\phantomsection\label{result-lemma-oa-2}{\textbf{Lemma OA.2
(analytical).}} \emph{For prior \(\pi\in(0,1)\), let
\(\mu_-=\operatorname{logistic}(\operatorname{logit}\pi-2/b)\),
\(\mu_+=\operatorname{logistic}(\operatorname{logit}\pi+2/b)\), and
\(\omega=\min\{\mu_-,1-\mu_+\}\). Replace (A1)--(A3) by
\(c_L<B_{r_1}(\mu_-)\), \(B_{r_0}(\pi)<c_H<B_{r_1}(\mu_+)\), and
\(\Delta_T(r_0)<k<(1-1/b)\rho\omega\Delta_T(r_1)\). Then the benchmark
entry reversal and unique trading outcomes hold.}

Prior odds multiply the order-flow likelihood ratio. The uniform
residual lower bound is therefore \(\rho\omega\Delta_T\), while the
upper bound is still \(\Delta_T\). Every inference and global deviation
argument follows with these bounds. Under full Laplace orders,
\(x^*=b[\operatorname{logit}\tau-\operatorname{logit}\pi]/2\) lies in
\((0,1)\), and total entry is
\(\rho+(1-\rho)[\pi\alpha_H+(1-\pi)\alpha_L]>\rho\). Strict feasibility
at the equal prior persists for nearby priors. Neither supplemental
lemma is required to generate a reported numerical value in the present
manuscript.

\subsection{A.7. Full proof and probability formulas for Proposition
A.7}\label{oa-a-signals}

The auction, value supports, order interval, noise law, and
preparation-cost distribution are those in A.1. Only the information
observations change. Let \(a,d\in(1/2,1)\) denote investor and buyer
signal accuracy. Define

\[
\begin{aligned}
\mu_-&=(1-a)+(2a-1)m,&\mu_+&=(1-a)+(2a-1)M,\\
\phi_+(\mu)&=\frac{d\mu}{d\mu+(1-d)(1-\mu)},&
\phi_-(\mu)&=\frac{(1-d)\mu}{(1-d)\mu+d(1-\mu)},\\
w_H(r)&=t_H(r)-t_0(r),&w_L(r)&=t_L(r)-t_0(r).
\end{aligned}
\tag{OA.28}\label{eq:oa-oa-28}
\]

The complete sufficient conditions for Proposition A.7 are

\[
\begin{aligned}
&c_L<B_{r_1}(\phi_-(\mu_-)),\\
&B_{r_0}(d)<c_H<B_{r_1}(\phi_+(\mu_+)),\\
&(2a-1)\{\Delta_T(r_0)+(1-\rho)(2d-1)w_H(r_0)\}<k\\
&\hspace{25mm}<(1-1/b)m(2a-1)\rho\Delta_T(r_1).
\end{aligned}
\tag{OA.29}\label{eq:oa-oa-29}
\]

I prove unique zero orders and entry \(\rho\) at \(r_0\), unique maximum
correctly signed signal-contingent orders at \(r_1\), and strictly
higher entry at \(r_1\). The result allows \(d>a\) and arbitrary mixed
orders conditional on the investor's signal.

I use \(T\) for the trader's signal, \(Y\) for the buyer's private
signal, and \(\Theta\) for fundamental quality. Conditional on \(H\),
the probabilities of \(T=+\) and \(Y=+\) are \(a\) and \(d\);
conditional on \(L\) they are \(1-a\) and \(1-d\). Conditional
independence gives the joint signal law as the product of these
probabilities within each fundamental state. Equal fundamental priors
imply equal marginal probabilities for each trader signal.

For arbitrary signal-contingent mixed orders, let
\(a_+(x)=\int f(x-q)d\sigma_+(q)\) and
\(a_-(x)=\int f(x-q)d\sigma_-(q)\). Then

\[
\lambda_X=\frac{a_+}{a_++a_-}\in[m,M],\qquad
\mu_X=\frac{a a_++(1-a)a_-}{a_++a_-}=(1-a)+(2a-1)\lambda_X.
\tag{OA.30}\label{eq:oa-oa-30}
\]

Since \(X\) depends on quality only through \(T\), \(\Theta\) and \(X\)
are conditionally independent given \(T\). Also \(Y\) is independent of
\(X\) conditional on \(\Theta\). The latter remains true after applying
a measurable price function, and implies

\[
\Pr(H\mid P,Y=y)=\phi_y(\mu_P).
\tag{OA.31}\label{eq:oa-oa-31}
\]

The maps \(\phi_+,\phi_-\) are continuous and strictly increasing on
\((0,1)\); the favorable private signal gives the larger posterior. The
public posterior based on price inherits \([\mu_-,\mu_+]\) by
conditional expectation. The lowest possible joint posterior is
consequently \(\phi_-(\mu_-)\). The low-cost restriction in
\eqref{eq:oa-oa-29} and declining gross profits in \(r\) make low-cost
preparation optimal in both economies under any candidate equilibrium.

Define high-cost indicators
\(I_y(P)=\mathbf1\{B_r(\phi_y(\mu_P))\ge c_H\}\). They obey
\(I_+\ge I_-\). Since \(Y\)'s law conditional on quality does not change
with \(X\) or the trader's signal, true-state entry conditional on \(P\)
is

\[
\begin{aligned}
e_H(P)&=\rho+(1-\rho)\{d I_+(P)+(1-d)I_-(P)\},\\
e_L(P)&=\rho+(1-\rho)\{(1-d)I_+(P)+dI_-(P)\}.
\end{aligned}
\tag{OA.32}\label{eq:oa-oa-32}
\] In particular,

\[
e_H-e_L=(1-\rho)(2d-1)(I_+-I_-),\qquad
e_H\ge e_L\ge\rho.
\tag{OA.33}\label{eq:oa-oa-33}
\]

Let \(D=e_Hw_H-e_Lw_L\). As \(w_H-w_L=\Delta_T\) and \(w_L>0\),

\[
D=e_L\Delta_T+(e_H-e_L)w_H
\in[\rho\Delta_T,\ \Delta_T+(1-\rho)(2d-1)w_H].
\tag{OA.34}\label{eq:oa-oa-34}
\]

Conditional pricing is \(P=t_0+e_L(P)w_L+\mu_XD(P)\). Both entry
probabilities are functions of price, even before proving sufficiency.
The positive coefficient \(D(P)\) therefore makes \(\mu_X\) measurable
with respect to price by the explicit inversion. The tower property
gives \(\mu_P=\mu_X\) and hence also reveals \(\lambda_X\) through the
affine inverse in \eqref{eq:oa-oa-30}.

For construction, set the entry indicators using \(\phi_y(\mu)\) and
write

\[
P(\mu)=t_0+\mu e_H(\mu)w_H+(1-\mu)e_L(\mu)w_L.
\tag{OA.35}\label{eq:oa-oa-35}
\]

Both entry probabilities are nondecreasing in \(\mu\). For
\(\mu_2>\mu_1\), decompose its change into a belief change at the old
entry probabilities and an entry change evaluated at the new belief:

\[
\begin{aligned}
P(\mu_2)-P(\mu_1)
={}&(\mu_2-\mu_1)D(\mu_1)
+\mu_2[e_H(\mu_2)-e_H(\mu_1)]w_H\\
&+(1-\mu_2)[e_L(\mu_2)-e_L(\mu_1)]w_L>0.
\end{aligned}
\tag{OA.36}\label{eq:oa-oa-36}
\]

Thus the same measurable-inverse construction works even when the entry
indicators jump simultaneously or one never changes. This argument does
not rely on differentiating an indicator.

Conditional on \(T=+\), the trader assigns quality probability \(a\)
even after its own randomized order and the resulting flow are
specified. Conditional on \(T=-\), it assigns \(1-a\). This remains
valid under unilateral orders because the deviation introduces no
additional quality information conditional on the signal. Its expected
terminal target payoff at a flow \(x\) is therefore
\(t_0+e_L(P(x))w_L+aD(P(x))\) for \(T=+\), and the corresponding
expression with \(1-a\) for \(T=-\). Subtracting \eqref{eq:oa-oa-35} and
using \eqref{eq:oa-oa-30} gives

\[
A_+=(2a-1)(1-\lambda_X)D,\qquad
A_-=(2a-1)\lambda_XD.
\tag{OA.37}\label{eq:oa-oa-37}
\]

The residuals have strict signs and obey

\[
m(2a-1)\rho\Delta_T\le A_+,A_-
\le(2a-1)\{\Delta_T+(1-\rho)(2d-1)w_H\}=: \overline A.
\tag{OA.38}\label{eq:oa-oa-38}
\]

The weak restriction in \eqref{eq:oa-oa-29} makes any correctly signed
order of magnitude \(s>0\) earn at most \(s(\overline A-k)<0\). Hence
zero is necessary under all candidate equilibria. With zero orders the
public posterior is the prior, the buyer's favorable posterior is \(d\),
and the expensive-entry exclusion in \eqref{eq:oa-oa-29} gives entry
\(\rho\). Constant competitive pricing constructs the weak equilibrium.

At high strength, the convolution regularity applies to either bounded
signal residual and gives \(U'(s)\ge(1-1/b)m(2a-1)\rho\Delta_T-k>0\).
Thus the necessary and unique trading profile is full correctly signed
orders by trader signal. Bayes' rule, \eqref{eq:oa-oa-35}, and optimal
preparation construct the equilibrium. The strict high-cost upper
inequality ensures that \(Y=+\) and public beliefs sufficiently near
\(\mu_+\) induce preparation. Full Laplace orders attain those beliefs
on a positive-probability event; the conditional probability of \(Y=+\)
is positive. Total entry strictly exceeds \(\rho\). The proof allows
\(d>a\) throughout.

For numerical implementation with either private signal, let
\(\tau=(c_H-g_L)/(g_H-g_L)\) and solve \(\phi_y(\mu)=\tau\):

\[
\begin{aligned}
\mu_{\mathrm{req},+}&=\frac{\tau(1-d)}{d(1-\tau)+\tau(1-d)},\\
\mu_{\mathrm{req},-}&=\frac{\tau d}{(1-d)(1-\tau)+\tau d},\\
\lambda_{\mathrm{req},y}&=\frac{\mu_{\mathrm{req},y}-(1-a)}{2a-1}.
\end{aligned}
\tag{OA.39}\label{eq:oa-oa-39}
\]

These formulas presume an interior fundamental threshold. Thresholds
outside \([0,1]\) imply always or never entry and are handled directly
by the profit comparison. For a requirement strictly between \(m\) and
\(M\), full Laplace orders give
\(x_y^*=b\operatorname{logit}(\lambda_{\mathrm{req},y})/2\).
Requirements below the lower bound or above the upper bound imply always
or never entry; equality at a plateau follows the tie rule. The buyer's
unfavorable signal need not be assumed to preclude high-cost entry in a
parameter sweep.

The fundamental-conditioned flow densities are

\[
f_H^X(x)=a f(x-1)+(1-a)f(x+1),\qquad
f_L^X(x)=(1-a)f(x-1)+a f(x+1).
\tag{OA.40}\label{eq:oa-oa-40}
\]

Let \(Q_{\theta,y}=\Pr(X\ge x_y^*\mid\Theta=\theta)\), with the
always/never cases evaluated accordingly. Conditional entry is

\[
\begin{aligned}
\bar e_H&=\rho+(1-\rho)[dQ_{H,+}+(1-d)Q_{H,-}],\\
\bar e_L&=\rho+(1-\rho)[(1-d)Q_{L,+}+dQ_{L,-}].
\end{aligned}
\tag{OA.41}\label{eq:oa-oa-41}
\]

Thus \(\mathsf E=(\bar e_H+\bar e_L)/2\), \(\mathsf O_H=\bar e_H/2\),
and \(\mathcal R_T=t_0+(\bar e_Hw_H+\bar e_Lw_L)/2\). Independently
integrating the joint law of \((\Theta,T,Y,Z,C)\) must reproduce these
expressions. Using the marginal law of \(Y\) instead of its conditional
law given quality would give the wrong entry and pricing objects.

\subsection{A.8. Verifiable-value bargaining}\label{oa-a-bargaining}

The bargaining comparison uses a different acquisition institution,
specified here. After diligence, all buyer values are verifiable. The
seller can implement a sale to the runner-up at that buyer's value, with
acceptance at zero buyer surplus; therefore the disagreement payoff in
negotiating with the best buyer is the next-best value \(z\). The best
buyer has value \(V\ge z\) and disagreement payoff zero. For
\(0<\eta<1\), its Nash transfer solves

\[
\max_{P\in[z,V]}(P-z)^\eta(V-P)^{1-\eta}.
\tag{OA.42}\label{eq:oa-oa-42}
\]

For \(V>z\), strict concavity of the log objective gives the unique
transfer \(P=(1-\eta)z+\eta V\). When \(V=z\), feasible transfers
collapse to that common value. At \(\eta=0\) I use the continuous
limiting transfer \(z\). With no challenger, the same rule has fallback
zero and pays the seller \(\eta R\). Thus no-entry proceeds are
specified by the institution as well as entry proceeds.

The highest-value buyer receives the target. For \(R\le\ell\), the high
and low target payments are \((1-\eta)R+\eta h\) and
\((1-\eta)R+\eta\ell\). For \(R>\ell\), they are \((1-\eta)R+\eta h\)
and \((1-\eta)\ell+\eta R\). Their difference is

\[
\eta(h-\ell)+(1-2\eta)(R-\ell)_+.
\tag{OA.43}\label{eq:oa-oa-43}
\]

A challenger receives \((1-\eta)(\theta-R)\) when \(\theta\ge R\) and
zero otherwise. Taking expectations gives Proposition A.8 and applying
FOSD yields its sign comparisons. For a uniform incumbent,

\[
\begin{aligned}
\Delta_\eta&=\eta(h-\ell)+(1-2\eta)\frac{(r-\ell)^2}{2r},\\
G_{H,\eta}&=(1-\eta)(h-r/2),\\
G_{L,\eta}&=(1-\eta)\frac{\ell^2}{2r}.
\end{aligned}
\tag{OA.44}\label{eq:oa-oa-44}
\]

The transfer is always feasible between fallback and winning value. For
\(\eta>1/2\) the competition derivative of the spread is negative; the
spread itself remains positive because the high challenger has the
higher value. The endpoint \(\eta=1\) extinguishes buyer rents and is
not used to infer a positive-cost entry equilibrium. The comparison does
not solve a private-value bargaining or first-price auction game;
verifiability and fallback enforceability are part of the stated
institution.

\subsection{A.9. Matched welfare and the external-dividend
diagnostic}\label{oa-a-welfare}

I couple the feedback and price-hidden economies at \(r=r_1\) using the
same \((\Theta,R,C,Z)\) and full informed orders. The strong trading
bound applies in the price-hidden economy with entry probability
\(\rho\), so this coupling uses equilibrium behavior in both
environments. The high-cost type remains out when price information is
unavailable because \(B_{r_1}(1/2)<c_H\); all low-cost realizations
enter in both environments.

Without challenger entry, allocation value is \(R\mathbf1\{R\ge p\}\).
With entry it is \(\max(R,\theta)\) because \(\theta>p\). Splitting at
\(R=p\) and \(R=\theta\) gives

\[
\max(R,\theta)-R\mathbf1\{R\ge p\}
=(\theta-\max(p,R))_++p\mathbf1\{R<p\}.
\tag{OA.45}\label{eq:oa-oa-45}
\]

If \(R<p\), both sides equal \(\theta\); if \(R\ge p\), both sides equal
\((\theta-R)_+\). Conditional on the signal and preparation cost, \(R\)
retains its independent uniform law, so incremental expected allocation
value net of cost is

\[
B_r(\mu_X)-C+p\Pr(R<p)=B_r(\mu_X)-C+\frac{p^2}{r}.
\tag{OA.46}\label{eq:oa-oa-46}
\]

Additional entry occurs only when its first two terms sum to a
nonnegative value. The final term is strictly positive. The statement is
about conditional expectations: each additional preparation decision
contributes at least \(p^2/r\) in conditional expectation, given price
information and preparation cost. A realized low-value preparation need
not generate a positive realized net gain, since \((\ell-\max(p,R))_+\)
can fall short of the cost paid. For atomic costs, integrate against the
marginal full-order flow density \(g(x)=[f(x-1)+f(x+1)]/2\) to obtain

\[
\Delta\mathcal W=(1-\rho)\int g(x)
\left[B_r(\mu_X(x))-c_H+\frac{p^2}{r}\right]
\mathbf1\{B_r(\mu_X(x))\ge c_H\}\,dx>0.
\tag{OA.47}\label{eq:oa-oa-47}
\]

For an atomless high-cost component \(H_H\), replace the integrand by

\[
(1-\rho)\int\left[B_r(\mu_X(x))-c+\frac{p^2}{r}\right]
\mathbf1\{c\le B_r(\mu_X(x))\}\,H_H(dc).
\tag{OA.48}\label{eq:oa-oa-48}
\]

The strict-support conditions ensure positive additional entry. Fubini
is valid because values and the specified cost supports are bounded. The
realized trading cost is the same in both environments; any resource
interpretation of that cost therefore cancels. Transfers between
bidders, shareholders, market makers, and noise traders are not
additional allocation surplus.

Target proceeds increase by averaging \(t_\theta-t_0>0\) over each
state's additional-entry probability. For atomic costs the difference is
\((1-\rho)\{\alpha_H(t_H-t_0)+\alpha_L(t_L-t_0)\}/2>0\); for atomless
costs integrate the corresponding state-specific tail probabilities over
the high-cost component. As an independent numerical check, compute
total expected allocation minus the full expected preparation cost in
each environment:

\[
\mathcal W=\frac{r^2-p^2}{2r}
+\frac12\sum_\theta\bar e_\theta\left(g_\theta+\frac{p^2}{r}\right)
-\mathbb E[C\mathbf1\{\text{entry}\}],
\tag{OA.49}\label{eq:oa-oa-49}
\]

valid in the benchmark support region. Subtracting the price-hidden
value must agree with \eqref{eq:oa-oa-47} or \eqref{eq:oa-oa-48}.

For the matched-level diagnostic, attach a deterministic payoff
\(D_0=\mathcal R_T^{\mathrm{feedback}}-\mathcal R_T^{\mathrm{hidden}}\)
to the traded claim in the price-hidden economy. Its competitive price
shifts by \(D_0\) and \(V_T-P\) is unchanged. The buyer's acquisition
profits and information set do not change. This matches mean prices
without manufacturing information. The external payoff is not counted as
a real surplus gain and is not an available reserve or acquisition
payment.

\subsection{A.10. Continuous acquisition values and sale-design
regularity}\label{oa-a-design}

For realized challenger value \(v\), define pointwise seller revenue

\[
T(R,v,p)=
\begin{cases}
p\mathbf1\{R\ge p\},&v<p,\\
\max\{p,\min(R,v)\},&v\ge p,
\end{cases}
\qquad
G(R,v,p)=(v-\max(p,R))_+.
\tag{OA.50}\label{eq:oa-oa-50}
\]

These formulas include no sale when both values fail the reserve. They
are bounded, Borel, and can be integrated first over \(R\) and then over
any conditional challenger distribution. The latter order is convenient
when a reserve cuts through a value band. At a value atom equal to the
reserve, the bid meets the reserve even when its acquisition rent is
zero; the cost of preparation was paid before the realized value was
learned.

For uniform \(R\) and a fixed \(v\), a useful complete kernel is

\[
\begin{array}{ll}
v<p:&t_v=t_0,\quad g_v=0,\\
p\le v<r:&t_v=v-\dfrac{v^2-p^2}{2r},\quad g_v=\dfrac{(v-p)^2}{2r}+\dfrac{p(v-p)}r,\\
p\le r\le v:&t_v=\dfrac r2+\dfrac{p^2}{2r},\quad g_v=v-t_v,\\
r<p\le v:&t_v=p,\quad g_v=v-p,
\end{array}
\tag{OA.51}\label{eq:oa-oa-51}
\]

with \(t_0=p(1-p/r)\) for \(p\le r\) and \(t_0=0\) for \(p>r\). The
third line presumes \(p\le r\). The kernels agree at \(v=r\) and
\(p=r\). At \(p=v\), the challenger remains admissible at zero
acquisition rent; moving the reserve above a value atom can therefore
cause a jump in target proceeds. In the second line
\(g_v=(v^2-p^2)/(2r)\). Direct integration of \eqref{eq:oa-oa-50},
rather than the shortened kernels alone, is the independent numerical
oracle.

For the atomless-value diagnostic, a binary class \(S\) is equally
likely. Conditional on \(S=L\), \(V\) is uniform on
\([\ell-\varepsilon_V,\ell+\varepsilon_V]\); conditional on \(S=H\), it
is uniform on \([h-\varepsilon_V,h+\varepsilon_V]\). The investor
observes \(S\), not the exact \(V\). Preparation reveals \(V\). The
independent incumbent remains uniform. Conditional class values have
equal width but different means. This is a binary-information economy
with atomless acquisition values, not a perfectly informed trader
observing a continuous value.

When \(p<\ell-\varepsilon_V<\ell+\varepsilon_V<r<h-\varepsilon_V\),

\[
\begin{aligned}
t_L&=\ell-\frac{\ell^2+\varepsilon_V^2/3-p^2}{2r},\qquad
g_L=\frac{\ell^2+\varepsilon_V^2/3-p^2}{2r},\\
t_H&=\frac r2+\frac{p^2}{2r},\qquad g_H=h-t_H,\qquad
\Delta_T=\frac{(r-\ell)^2}{2r}+\frac{\varepsilon_V^2}{6r}.
\end{aligned}
\tag{OA.52}\label{eq:oa-oa-52}
\]

When the reserve is above the entire low band but below \(r\),
\(t_L=t_0\), \(g_L=0\), and the high-class formulas are unchanged.
Conditional on class and flow, exact \(V\) retains its within-class
distribution, so the benchmark price and residual proof applies to these
class-averaged payoffs whenever \(\Delta_T>0\) and a positive entry
floor holds. In this region \(t_L=t_0\) is permitted: the price
construction remains strictly increasing on the bounded posterior
interval because \(\Delta_T\mu>0\).

For a general reserve, the object the seller faces after choosing \(p\)
is a complete continuation

\[
\sigma=(\sigma_H,\sigma_L;\ P(\cdot);\ \mu_P;\ \text{preparation rule})\in\mathcal E(p,r),
\tag{OA.53}\label{eq:oa-oa-53}
\]

the conditional order distributions, the price mapping on flow, the
buyer's beliefs at every realized price including price atoms, and the
preparation rule at every price and cost. Its beliefs are the public
price posterior, not necessarily the raw order-flow posterior. If entry
vanishes on a set of flows, the inversion denominator vanishes there,
those flows share the no-entry price, and the buyer conditions on the
whole preimage of that price. The construction used in C.6, which prices
every positive-entry flow by the benchmark inversion and pools exactly
the zero-entry preimage at \(t_0\), is one admissible member of
\(\mathcal E(p,r)\) when it validates. It is not the only one.
\hyperref[oa-a-pools]{Section A.11} exhibits, in the reserve game
itself, a one-parameter family of continuations with the same orders and
different pools, each of which validates. Two continuations with the
same orders but different price mappings are different elements of
\(\mathcal E(p,r)\), and the numerical record identifies a continuation
by the whole object \eqref{eq:oa-oa-53}. If \(\Delta_T=0\) or all entry
is impossible, the informative-price argument cannot be imported.
\hyperref[oa-c-reserve]{Section C.6} therefore requires direct pricing,
posterior, and entry verification at such reserves.

For the local decomposition in the paper, impose a neighborhood \(I\)
contained strictly in \(0<p<\ell\), a pure strategy branch
\(q_\theta(p)\) continuously differentiable on \(I\), and an atomless
preparation-cost CDF \(H_C\) continuously differentiable on the compact
profit range. Suppose \(\mu_\theta(z;p)\) is differentiable in \(p\) for
almost every \(z\) and there is an integrable function \(J(z)\)
dominating

\[
f(z)\left|h_C(B_{p,r}(\mu_\theta))\left[-p/r+(g_H-g_L)\partial_p\mu_\theta\right]\right|
\tag{OA.54}\label{eq:oa-oa-54}
\]

uniformly on a compact neighborhood of the reserve. These assumptions
are sufficient, not an assertion that every continuation has them. For
positive smooth noise with bounded log-density derivative and a bounded
derivative of each pure order, posterior derivatives can be bounded
directly by likelihood-ratio differentiation; bounded \(h_C\) then
supplies domination by a constant times \(f\). For Laplace noise, the
same calculation holds away from the finite shifted kinks, a null set in
\(z\), and the bounded log slope again supplies domination.

For \(\bar e_\theta(p)=\int f(z)H_C(B_{p,r}(\mu_\theta(z;p)))\,dz\),
dominated differentiation gives

\[
\frac{d\bar e_\theta}{dp}
=\int f(z)h_C(B_{p,r}(\mu_\theta))
\left[-\frac p r+(g_H-g_L)\frac{\partial\mu_\theta}{\partial p}\right]dz.
\tag{OA.55}\label{eq:oa-oa-55}
\]

With \(\mathsf E=(\bar e_H+\bar e_L)/2\) and
\(\mathcal R_T=t_0+\sum_\theta\bar e_\theta(t_\theta-t_0)/2\), the
product rule gives

\[
\frac{d\mathcal R_T}{dp}
=(1-\mathsf E)\left(1-\frac{2p}{r}\right)+\mathsf E\frac p r
+\frac12\sum_\theta\frac{d\bar e_\theta}{dp}(t_\theta-t_0).
\tag{OA.56}\label{eq:oa-oa-56}
\]

The calculation uses \(\partial_pt_0=1-2p/r\),
\(\partial_pt_H=\partial_pt_L=p/r\), and \(\partial_pB|_\mu=-p/r\). At
an atomic preparation threshold or a nonregular continuation, this
differentiation argument is unavailable. The level objective and full
conditional laws remain the correct objects.

A seller equilibrium requires a feasible continuation selection after
each reserve and a maximizing reserve under that selection. The
numerical envelope of found continuations is an exploration of this
problem, not a proof of existence, measurable selection, global
attainment, or optimality. These are the explicit open parts of the
sale-design result.

\subsection{A.11. Price pooling when preparation can
vanish}\label{oa-a-pools}

When preparation can vanish, identical investor orders can support
different price pools and different preparation outcomes. This
subsection proves the explicit family in Proposition A.10. Dow et al.
(\citeproc{ref-DowGoldsteinGuembel2017}{2017}) discuss continuations
distinguished by their pooling regions following their Lemma 1. The
reserve-game example makes the price rule part of the continuation that
the seller must face.

\textbf{Inputs and institution.} Take the benchmark declaration of C.0
with \(r=r_0\) and change only the reserve, to the value declared as
\texttt{pool\_reserve} in the registry:

\begin{Shaded}
\begin{Highlighting}[]
\NormalTok{h = 10; ell = 1; r = 1.2; p = 7;}
\NormalTok{rho = 0.25; c\_L = 1; c\_H = 6;}
\NormalTok{b = 2; k = 0.02; prior\_H = 0.5;}
\NormalTok{q\_H = 1; q\_L = {-}1;}
\NormalTok{entry\_at\_indifference = yes;}
\NormalTok{bid\_equal\_reserve\_is\_admissible = yes.}
\end{Highlighting}
\end{Shaded}

The reserve exceeds incumbent support and excludes the low-value
challenger, so this is the expanded reserve domain of A.10, not the
benchmark support \(p<\ell\). From \eqref{eq:oa-oa-51} with \(r<p\le v\)
for \(v=h\) and \(v<p\) for \(v=\ell\),

\[
t_0=t_L=g_L=0,\qquad t_H=7,\qquad g_H=3,\qquad \Delta_T=7.
\]

A payoff routine written for \(p<\ell\) must not be applied here; C.6b
evaluates these objects in exact rational arithmetic and by direct
integration of \eqref{eq:oa-oa-50}.

\textbf{Densities and the price family.} With \(b=2\) and full orders,

\[
\begin{aligned}
f(z)&=\frac14e^{-|z|/2},\qquad a_H(x)=f(x-1),\qquad a_L(x)=f(x+1),\\
\mu_X(x)&=\left[1+\exp\left\{-\frac{|x+1|-|x-1|}{2}\right\}\right]^{-1},
\end{aligned}
\tag{OA.57}\label{eq:oa-oa-57}
\]

which is \eqref{eq:oa-oa-16} at \(b=2\): the plateau \(m=(1+e)^{-1}\) on
\(x\le-1\), \(\operatorname{logistic}(x)\) on \((-1,1)\), and \(M=1-m\)
on \(x\ge1\). For each cutoff \(c\in[-\log2,0]\) define the price rule

\[
P_c(x)=
\begin{cases}
0,&x<c,\\
\tfrac74\mu_X(x),&x\ge c,
\end{cases}
\tag{OA.58}\label{eq:oa-oa-58}
\]

with the preparation policy: at price zero neither cost type prepares;
at every positive price the low-cost type prepares and the high-cost
type does not. The upper plateau \(x\ge1\) maps to the price
\(\tfrac74M\), a positive-price atom with probability
\(\tfrac12[S_Z(0)+S_Z(2)]\); the buyer conditions on it as an atom, and
its posterior is \(M\) because the plateau is a level set of \(\mu_X\).

\textbf{Noise CDF versus survival.} I write \(F_Z\) for the noise CDF
and \(S_Z=1-F_Z\) for its survival function, and keep the two apart
throughout, because the pooled posterior uses the CDF and the outcome
quantities use the survival function:

\[
F_Z(z)=
\begin{cases}
\tfrac12e^{z/2},&z\le0,\\
1-\tfrac12e^{-z/2},&z>0,
\end{cases}
\qquad
S_Z(z)=1-F_Z(z).
\tag{OA.59}\label{eq:oa-oa-59}
\]

Equation \eqref{eq:oa-oa-17} is \(S_Z\) at general \(b\).

\textbf{Buyer at price zero.} The zero price is one atom whose preimage
is the whole pool \(\{X<c\}\). Its posterior follows \eqref{eq:oa-oa-79}
with that preimage:

\[
\bar\mu_c=\Pr(H\mid X<c)=\frac{\int_{-\infty}^ca_H}{\int_{-\infty}^ca_H+\int_{-\infty}^ca_L}
=\frac{F_Z(c-1)}{F_Z(c-1)+F_Z(c+1)}.
\tag{OA.60}\label{eq:oa-oa-60}
\]

The likelihood ratio \(a_H/a_L\) is nondecreasing, so \(\bar\mu_c\) is
nondecreasing in \(c\) and, over the declared family,
\(\bar\mu_c\le\bar\mu_0=\tfrac12e^{-1/2}<\tfrac13\). Gross profit at the
pool is \(B(\bar\mu_c)=g_H\bar\mu_c=3\bar\mu_c<1=c_L\), so not preparing
is strictly optimal at price zero for both cost types. The point that
matters for the numerical layer is the case \(c=0\). There the raw flow
posterior \(\mu_X(x)\) exceeds \(\tfrac13\) on \((-\log2,0)\), and a
buyer who observed those flows one by one would prepare at low cost.
Those flows are not individually observed. They all produce price zero,
the buyer's information set at price zero is the pool, and the pooled
posterior governs. A candidate that lets the buyer act on \(\mu_X\)
inside the pool is not measurable with respect to the price and is
rejected for that reason, not for a numerical breach. C.6b carries this
as a negative control.

\textbf{Buyer at positive prices.} On \(x\ge c\) the price
\(\tfrac74\mu_X(x)\) is strictly increasing in \(\mu_X\), so the buyer
recovers \(\mu_X\) from the price on \((-1,1)\) and the plateau atom on
\(x\ge1\) carries posterior \(M\). Since \(c\ge-\log2\) and \(\mu_X\) is
nondecreasing, \(\mu_X(x)\ge\mu_X(-\log2)=\tfrac13\) on the
positive-price region, so \(3\mu_X\ge c_L\) and low-cost preparation is
optimal under the equality convention, with equality at \(c=-\log2\)
exactly at the cutoff. High-cost preparation is impossible at any price:
gross profit is at most \(g_H=3<6=c_H\). The zero-price and
positive-price regions do not collide, since
\(\tfrac74\mu_X\ge\tfrac7{12}>0\) on \(x\ge c\).

\textbf{Conditional pricing.} For \(x<c\) nobody is prepared, neither
the incumbent nor a low-value challenger can meet the reserve, and the
terminal value is zero. For \(x\ge c\) preparation occurs with
probability \(\rho\), and a sale, at exactly \(p\), occurs only when the
prepared challenger has value \(h\):

\[
\mathbb E[V_T\mid X=x]=\rho\,p\,\mu_X(x)=\tfrac74\mu_X(x)\quad(x\ge c),\qquad
\mathbb E[V_T\mid X=x]=0\quad(x<c).
\tag{OA.61}\label{eq:oa-oa-61}
\]

So \(P_c\) is the competitive price pointwise in flow on both regions.
This is the test. Average price equal to average payoff is implied by
it, and a wrong pointwise schedule can pass the average check, so C.6b
verifies \eqref{eq:oa-oa-61} on a flow mesh that includes the cutoff
from both sides and every breakpoint, and treats the average identity as
a separate necessary check.

\textbf{Global investor validation.} Against a fixed member of the
family the residuals of A.3, after the preparation response, are

\[
A_H(x)=\rho p[1-\mu_X(x)]\mathbf1\{x\ge c\},\qquad
A_L(x)=\rho p\,\mu_X(x)\mathbf1\{x\ge c\}.
\tag{OA.62}\label{eq:oa-oa-62}
\]

They vanish inside the pool, where the price carries no information and
no entry follows, and are nonnegative elsewhere; wrong-signed orders are
therefore dominated by zero. Let \(J_c=F_H(1)=F_L(1)\) as in
\eqref{eq:oa-oa-22}. The equality holds pointwise inside the integral by
Bayes' rule, since \(f(x-1)[1-\mu_X(x)]=f(x+1)\mu_X(x)\). The Laplace
density-ratio bound gives \(F_\theta(s)\ge e^{-(1-s)/b}J_c\) and
\(|F_\theta'(s)|\le F_\theta(s)/b\) by \eqref{eq:oa-oa-12}, hence by
\eqref{eq:oa-oa-13}

\[
U_\theta'(s)\ge\left(1-\frac1b\right)J_c-k
\]

for every \(s\in[0,1]\). Every cutoff in the family satisfies
\(c\le0<1\), so the tail \(x\ge1\) is active with entry \(\rho\) under
each member. On that tail \(1-\mu_X=m=(1+e)^{-1}>\tfrac14\), and its
probability in state \(H\) is \(S_Z(0)=\tfrac12\). Dropping the rest of
the active region,

\[
J_c\ge\frac{\rho p\,m}{2}>\frac{7}{32},
\qquad
U_\theta'(s)>\frac{7}{64}-\frac{1}{50}=\frac{143}{1600}>0.
\tag{OA.63}\label{eq:oa-oa-63}
\]

The rational bound \eqref{eq:oa-oa-63} is the proof, common to every
member of the family and to both types. A positive derivative found on a
finite grid is not the proof; it is an implementation check, and C.6b
keeps the two apart by writing the rational bound and the mesh minimum
in separate columns. Full correctly signed orders are globally optimal
against every price rule in the family. This is analytical existence for
the whole family. It says nothing about whether other continuations
exist in \(\mathcal E(7,r_0)\).

\Needspace{8\baselineskip}

\protect\phantomsection\label{result-proposition-oa-3}{\textbf{Proposition
OA.3 (price pooling at fixed orders; analytical).}} \emph{For every
\(c\in[-\log2,0]\), the price rule \eqref{eq:oa-oa-58} with the stated
preparation policy and full orders \((1,-1)\) is a continuation
equilibrium in \(\mathcal E(7,r_0)\): prices are competitive pointwise
in flow, beliefs at the zero-price atom are \eqref{eq:oa-oa-60}, the
buyer's preparation decision is optimal at every price information set,
and full correctly signed orders are globally optimal for both investor
types against the fixed price and preparation schedule. The orders are
the same in all members of the family; prices, beliefs, and preparation
differ.}

The proof is the four validations above. The statement does not
establish uniqueness of this family within \(\mathcal E(7,r_0)\), and it
is confined to lower-interval pools \(\{X<c\}\); other measurable pool
sets are not searched, and the numerical record says so through
\texttt{pricing\_family\_coverage} and
\texttt{exhaustive\_pricing\_search\ =\ false}. On the benchmark support
the family cannot arise: there \(c_L<B_r(m)\) gives entry at least
\(\rho\) at every flow and Proposition A.2 ties the price experiment to
the orders.

\textbf{Output quantities.} With \(S_Z\) from \eqref{eq:oa-oa-59},

\[
\begin{aligned}
\Pr(P=0)&=\frac{F_Z(c-1)+F_Z(c+1)}2,\\
\mathsf E_c&=\frac\rho2\{S_Z(c-1)+S_Z(c+1)\},\qquad
\Pr(\text{sale})=\frac\rho2S_Z(c-1),\\
\mathcal R_{T,c}&=\frac{\rho p}{2}S_Z(c-1),\qquad
\Pr(\text{two admissible bidders})=0.
\end{aligned}
\tag{OA.64}\label{eq:oa-oa-64}
\]

The admissible-challenger probability equals the sale probability,
because only a prepared high-value challenger can meet the reserve and
the incumbent never does. Both \(\mathsf E_c\) and \(\mathcal R_{T,c}\)
are strictly decreasing in \(c\). A higher cutoff pools more flows into
the uninformative price and suppresses preparation at flows that would
have justified it if observed separately. C.6b recomputes
\eqref{eq:oa-oa-64} in fifty-digit arithmetic and by independent
conditional integration, and reports the endpoints \(c=-\log2\) and
\(c=0\) through the registry keys \texttt{pool\_posterior\_cutoff\_*},
\texttt{pool\_entry\_cutoff\_*}, and \texttt{pool\_revenue\_cutoff\_*}.

\textbf{The seventeen-cutoff grid and the negative controls.} The
declared grid is \(c_j=-\log2\,(1-j/16)\) for \(j=0,\dots,16\), with the
endpoints included. Every grid member must pass the price, buyer, and
global-investor validations, and preparation and revenue must decrease
along the grid. That ordered check is a check of the implementation of
\eqref{eq:oa-oa-64}, not a continuum theorem. Three controls must fail,
each for its own economic reason. At \(c=-1\) some positive prices have
\(3\mu_X<1\), so the prescribed low-cost preparation is not optimal
there. At \(c=1\) the pooled zero-price posterior \(\bar\mu_1\) makes
low-cost preparation profitable, so the prescribed nonpreparation fails.
Inside the \(c=0\) pool, a buyer who acts on the raw posterior is
rejected as inconsistent with the observed-price information structure,
whatever its arithmetic. A fourth control is about identity rather than
validation: a deduplication key built from orders alone merges the two
valid endpoint continuations, which have the same orders and different
atoms, and the C.6b ledger must show that merge so that the key is never
used. Failed controls are stored as expected failures with their reason,
never as accepted rows.

\section{B. Computer-assisted equilibrium certificates}\label{oa-b}

The certificates in Proposition 3 establish an exact equilibrium inside
each reported bracket. Outward interval arithmetic encloses a
marginal-profit root for the low-value investor and verifies global
optimality for the high-value investor. The argument combines analytical
regularity and deviation bounds with finite interval calculations.

\subsection{B.1. Objects to be certified}\label{oa-b-objects}

Fix the benchmark primitives and a strength \(r\). Candidate orders are
\((q_H,q_L)=(1,-v)\) with \(v\in(0,1)\). Let
\(\tau=(c_H-g_L)/(g_H-g_L)\), \(c=(1-v)/2\), and

\[
\mu_v(x)=\operatorname{logistic}\left(\frac{|x+v|-|x-1|}{b}\right),\quad
x^*(r,v)=\frac{b\operatorname{logit}\tau+1-v}{2}.
\tag{OA.65}\label{eq:oa-oa-65}
\]

The certificate region has \(1/2<\tau<M_v\), where
\(M_v=\operatorname{logistic}((1+v)/b)\). Thus \(x^*\) is strictly
between \(-v\) and \(1\). Set \(e(x)=\rho+(1-\rho)\mathbf1\{x\ge x^*\}\)
and define \(A_H=e\Delta_T(1-\mu_v)\), \(A_L=e\Delta_T\mu_v\).

For a signed direction \(\epsilon\in\{-1,1\}\) and magnitude
\(s\in[0,1]\), the investor's per-unit convolution and marginal profit
are

\[
F_\epsilon(s)=\int f(x-\epsilon s)A_\epsilon(x)\,dx,
\qquad U_\epsilon'(s)=F_\epsilon(s)+sF_\epsilon'(s)-k,
\tag{OA.66}\label{eq:oa-oa-66}
\]

where \(A_+=A_H\) and \(A_-=A_L\). In the root condition
\(\Psi(r,v)=U_-'(v;1,-v)\), the derivative is with respect to the
deviating magnitude \(s\), holding the candidate schedule fixed. The
outer change of \(v\) recomputes that schedule. Confusing these
derivatives would solve the wrong equilibrium condition.

\subsection{B.2. Global low-type optimality using a Stieltjes
measure}\label{oa-b-concavity}

The Laplace location family under \(1>-v\) gives a nondecreasing
posterior, strictly increasing on \((-v,1)\). Because \(e(\mu)\) is
nondecreasing and bounded below by \(\rho\), the function \(A_L\) is
bounded, nondecreasing, and nonconstant. Its right-continuous
representative defines a finite nonzero positive measure \(dA_L\) by
\(dA_L((x,y])=A_L(y)-A_L(x)\).

Integration by parts on finite intervals, followed by passage to the
infinite endpoints, gives

\[
F_L'(s)=\int f'(x+s)A_L(x)\,dx=-\int f(x+s)\,dA_L(x).
\tag{OA.67}\label{eq:oa-oa-67}
\]

The boundary term vanishes because \(A_L\) is bounded and the Laplace
density tends to zero in both tails. The measure integral is finite
because \(f\) is bounded and \(dA_L\) has finite mass. It is strictly
positive before the minus sign because \(f>0\) and the measure is
nonzero. An entry jump contributes its positive atom to this measure.

The function \(f\) is globally Lipschitz. Therefore
\(s\mapsto\int f(x+s)dA_L(x)\) is Lipschitz, implying that \(F_L'\) is
absolutely continuous and differentiable almost everywhere. Except when
\(-s\) hits an atom of \(dA_L\), its derivative follows by dominated
differentiation. The exceptional set of such atoms is at most countable,
so

\[
F_L''(s)=-\int f'(x+s)\,dA_L(x),\qquad
|F_L''(s)|\le\frac1b\int f(x+s)\,dA_L(x)=-\frac1bF_L'(s)
\tag{OA.68}\label{eq:oa-oa-68}
\]

almost everywhere. Hence

\[
U_L''(s)=2F_L'(s)+sF_L''(s)
\le(2-s/b)F_L'(s)<0
\tag{OA.69}\label{eq:oa-oa-69}
\]

for \(b>1/2\) on the unit order interval. Since \(U_L'\) is absolutely
continuous, this strict almost-everywhere inequality makes it strictly
decreasing. Any interior root is therefore the unique global maximizer
of \(U_L\) over all correctly signed magnitudes, including the
endpoints.

\subsection{B.3. Continuity of the equilibrium root
equation}\label{oa-b-root}

Let \([v_-,v_+]\) lie strictly inside \((0,1)\) and preserve
\(1/2<\tau<M_v\). The threshold \(x^*(r,v)\) is continuous in \(v\). For
any convergent sequence of candidate magnitudes, the posterior and the
entry residuals converge pointwise except at the limiting entry
boundary. The density and its derivative are translated by bounded
amounts; they are dominated by constant multiples of \(f(x)\), with an
additional factor \(1/b\) for the derivative. The residuals are bounded
by \(\Delta_T\). Dominated convergence therefore gives continuity of
both \(F_L(v;v)\) and its unilateral derivative component, except for
individual density kink points that have zero Lebesgue mass and do not
change the integral. Thus \(\Psi(r,v)\) is continuous.

Outward enclosures satisfying

\[
\inf\Psi(r,v_-)>0,\qquad \sup\Psi(r,v_+)<0
\tag{OA.70}\label{eq:oa-oa-70}
\]

prove existence of an exact root \(v^*\in(v_-,v_+)\). The concavity
result proves optimality at that root against its own price schedule. It
does not prove that \(\Psi\) has only one root across all possible
candidate schedules; no such uniqueness is needed for Proposition 3.

\subsection{B.4. A global cover for favorable-information
deviations}\label{oa-b-cover}

The distributional second derivative of the Laplace density is

\[
f''=\frac{f}{b^2}-\frac{\delta_0}{b^2}.
\tag{OA.71}\label{eq:oa-oa-71}
\]

Away from zero, ordinary differentiation gives \(f''=f/b^2\). The first
derivative jumps from \(1/(2b^2)\) to \(-1/(2b^2)\) at zero,
contributing the atom \(-\delta_0/b^2\). This verifies
\eqref{eq:oa-oa-71}, for example by integrating twice against a
compactly supported smooth test function.

Convolving with a bounded measurable \(A_H\) yields
\(F_H''=(F_H-A_H)/b^2\) as distributions and almost everywhere as
functions. Since \(F_H,A_H\in[0,\Delta_T]\), the weak second derivative
is bounded by \(\Delta_T/b^2\) in absolute value. The first derivative
has an absolutely continuous, Lipschitz representative; combined with
\(|F_H'|\le\Delta_T/b\), this gives

\[
|U_H''(s)|\le L_U=\frac{2\Delta_T}{b}+\frac{\Delta_T}{b^2},\qquad s\in[0,1]\ \text{a.e.}
\tag{OA.72}\label{eq:oa-oa-72}
\]

For \(s_j=j/n\), evaluate interval enclosures of \(U_H'(s_j)\)
\textbf{uniformly for every \(v\) in the root bracket}. Every point of
the order interval lies within \(1/(2n)\) of a mesh point, so the global
bound is

\[
\Gamma_H:=\min_j\inf U_H'(s_j;[v_-,v_+])-\frac{L_U}{2n}.
\tag{OA.73}\label{eq:oa-oa-73}
\]

A strictly positive lower enclosure for \(\Gamma_H\) proves that buying
the maximum amount is uniquely optimal. Checking only the numerical root
midpoint does not suffice: the exact root is known only to lie inside
the bracket. The interval cover must hold for the whole bracket. All
wrong-signed orders are dominated by zero by the residual signs already
proved.

\subsection{B.5. Elementary antiderivatives and infinite
tails}\label{oa-b-integrals}

Split the integration domain at \(-v\), \(x^*\), \(1\), and the density
center \(\zeta=\epsilon s\). Within the posterior's central region let
\(t=e^{(x-c)/b}\), so \(\mu_v=t^2/(1+t^2)\) and \(dx=b\,dt/t\). I use
the following primitives, each verified by differentiation:

\[
\begin{array}{c|cc}
&\displaystyle\int e^{x/b}(\cdot)\,dx&\displaystyle\int e^{-x/b}(\cdot)\,dx\\[2pt]
1-\mu_v&b e^{c/b}\arctan t&b e^{-c/b}(-t^{-1}-\arctan t)\\[2pt]
\mu_v&b e^{c/b}(t-\arctan t)&b e^{-c/b}\arctan t
\end{array}
\tag{OA.74}\label{eq:oa-oa-74}
\]

For \(x<\zeta\), the density multiplier is \(e^{-\zeta/b}e^{x/b}/(2b)\);
for \(x>\zeta\) it is \(e^{\zeta/b}e^{-x/b}/(2b)\). Multiply the
relevant primitive by this constant, by \(e\Delta_T\), and evaluate it
at the segment endpoints. On each segment the marginal-density
multiplier for \(F_\epsilon'\) is
\(\epsilon\operatorname{sgn}(x-\zeta)/b\), so the same segment integral
gives the derivative. This avoids differentiating the candidate entry
threshold during a unilateral deviation.

Outside \([-v,1]\) the posterior is constant. For a segment
\([a_0,a_1]\) left of \(\zeta\), its density mass is
\([e^{(a_1-\zeta)/b}-e^{(a_0-\zeta)/b}]/2\); right of \(\zeta\) it is
\([e^{-(a_0-\zeta)/b}-e^{-(a_1-\zeta)/b}]/2\). The infinite endpoint
terms vanish analytically. Multiplying by the segment's constant
residual integrates both infinite tails without truncation.

Endpoint ordering must be certified. When the low trader is evaluated at
\(s=v\), the center is identically \(-v\) and must be coalesced
algebraically, not treated as two uncertain endpoints. The same applies
at \(s=1\) for the high trader. For distinct uncertain endpoints, prove
the ordering by disjoint interval bounds. If an interval center overlaps
an entry boundary or another cut, subdivide the parameter box or the
order cell until the ordering can be proved; do not select an order
using only midpoints.

\subsection{B.6. Certificate acceptance and what it
establishes}\label{oa-b-acceptance}

I use exact decimal input strings and interval operations rounded
outward. In an interval test, the infimum and supremum of a displayed
function evaluation mean the endpoints of its outward interval
enclosure; the underlying equilibrium function remains scalar-valued.
The starting precision and mesh are specified in Appendix C.2. An
enclosure must retain its full endpoint representation; a printed
floating-point midpoint is not the certificate. For a bracket to
establish the stated result, I require all of the following: the model
support inequalities, low-cost participation and high-cost prior
exclusion, strict threshold ordering throughout the bracket, opposite
enclosed endpoint signs in \eqref{eq:oa-oa-70}, positive \(\Gamma_H\) in
\eqref{eq:oa-oa-73}, and a positive pooling-existence margin at the same
strength. Entry is enclosed by evaluating the exact tail formula
throughout the bracket. Pairwise strictly ordered entry intervals
establish the cross-economy increase.

Failure of an enclosure to exclude zero is an unresolved certificate,
not evidence of a profitable deviation or nonexistence. A failed sign in
a claimed successful certificate is an acceptance failure. Raising
precision or refining a bracket is legitimate only with the complete
failed and successful record retained.

The result is computer-assisted existence for exact equilibria inside
specified intervals. It does not prove uniqueness of the informative
profile, exclude mixed equilibria, or certify an entire continuous
branch between nodes. A continuation curve between certified nodes is a
numerical diagnostic until additional interval and continuation
arguments establish it. This distinction governs the correspondence
figure.

\section{C. Numerical exercises and quantity registry}\label{oa-c}

\subsection{C.0. Input declarations, output conventions, and
acceptance}\label{oa-c-contract}

This section is the complete numerical contract. The literal values in
its input declarations specify experiments; the manuscript's
quantitative placeholder fields are filled only by validated outputs or
these declared inputs. Mathematical constants, equation and result
labels, citation metadata, and software identifiers are structural, not
estimated quantities. The empirical scaffold contributes no sample size
or empirical estimate to the registry.

All parameters below are exact decimal inputs, converted without a
binary-floating intermediate when interval certification is required. A
numerical result records its full parameter vector, noise law, cost law,
information structure, sale institution, tie convention, and branch.
Defaults must not silently cross from one exercise to another. Scalar
probabilities are stored as probabilities; percentage formatting is a
separate presentation step.

\Needspace{5\baselineskip}

\textbf{Input declaration: benchmark.}

\begin{Shaded}
\begin{Highlighting}[]
\NormalTok{h = 10; ell = 1; p = 0.5; rho = 0.25;}
\NormalTok{c\_L = 1; c\_H = 6; b = 2; k = 0.02;}
\NormalTok{r\_weak = 1.2; r\_strong = 3; r\_collapse = 3.6;}
\NormalTok{noise = Laplace; cost = atomic; fundamental\_prior\_H = 0.5;}
\NormalTok{entry\_at\_indifference = yes; bid\_equal\_reserve\_is\_admissible = yes;}
\NormalTok{cost\_halfwidth = 0.1; value\_band\_halfwidth = 0.05;}
\NormalTok{binary\_alternative\_reserve = 1.01; atomless\_alternative\_reserve = 1.1.}
\end{Highlighting}
\end{Shaded}

\Needspace{5\baselineskip}

\textbf{Input declaration: moderate acquisition values.}

\begin{Shaded}
\begin{Highlighting}[]
\NormalTok{h = 2; ell = 1; p = 0.5; rho = 0.25;}
\NormalTok{c\_L = 0.3; c\_H = 0.89; b = 2; k = 0.002;}
\NormalTok{r\_weak = 1.05; r\_strong = 1.5;}
\NormalTok{noise = Laplace; cost = atomic; fundamental\_prior\_H = 0.5.}
\end{Highlighting}
\end{Shaded}

\Needspace{5\baselineskip}

\textbf{Input declaration: complementary signals.}

\begin{Shaded}
\begin{Highlighting}[]
\NormalTok{h = 10; ell = 1; p = 0.5; rho = 0.85;}
\NormalTok{c\_L = 1; c\_H = 7.14; b = 2; k = 0.015;}
\NormalTok{r\_weak = 1.1; r\_strong = 2.3; a = 0.70; d = 0.75;}
\NormalTok{noise = Laplace; cost = atomic; fundamental\_prior\_H = 0.5;}
\NormalTok{T\_and\_Y\_independent\_given\_fundamental = yes.}
\end{Highlighting}
\end{Shaded}

\Needspace{5\baselineskip}

\textbf{Input declaration: numerical controls.}

\begin{Shaded}
\begin{Highlighting}[]
\NormalTok{quadrature\_absolute\_target = 1e{-}11;}
\NormalTok{quadrature\_relative\_target = 1e{-}11;}
\NormalTok{probability\_acceptance = 1e{-}8;}
\NormalTok{price\_identity\_acceptance = 1e{-}8;}
\NormalTok{entry\_optimality\_acceptance = 1e{-}8;}
\NormalTok{deviation\_gain\_acceptance = 1e{-}7;}
\NormalTok{independent\_formula\_acceptance = 1e{-}9;}
\NormalTok{initial\_order\_intervals = 400;}
\NormalTok{refined\_order\_intervals = 800;}
\NormalTok{initial\_flow\_halfwidth = 40;}
\NormalTok{interval\_decimal\_precision = 50;}
\NormalTok{certificate\_derivative\_intervals = 200;}
\NormalTok{certificate\_refined\_intervals = 400;}
\NormalTok{probability\_display\_decimals = 6;}
\NormalTok{parameter\_display = shortest\_exact\_decimal;}
\NormalTok{certificate\_display = outward\_interval;}
\end{Highlighting}
\end{Shaded}

\Needspace{5\baselineskip}

\textbf{Input declaration: continuation, event, and outcome controls.}

\begin{Shaded}
\begin{Highlighting}[]
\NormalTok{pool\_reserve = 7;}
\NormalTok{price\_pool\_cutoff\_grid = {-}log2 * (1 {-} j/16), j = 0..16;}
\NormalTok{price\_pool\_negative\_controls = \{c = {-}1, c = 1, raw posterior inside the c = 0 pool,}
\NormalTok{                                orders{-}only deduplication key\};}
\NormalTok{pricing\_family\_coverage = recorded per row (text);}
\NormalTok{exhaustive\_pricing\_search = false unless a separate exhaustive argument is recorded;}
\NormalTok{order\_identity\_tolerance = 1e{-}6;}
\NormalTok{atom\_identity\_tolerance = 1e{-}8;}
\NormalTok{event\_arithmetic\_digits = 50;}
\NormalTok{event\_sign\_resolution = 1e{-}40;}
\NormalTok{event\_offsets = \{1e{-}4, 1e{-}6, 1e{-}8\};}
\NormalTok{tie\_rule = entry\_at\_indifference; bid\_equal\_reserve\_admissible;}
\NormalTok{           upper plateau enters at tau = M; lower plateau enters at B(m) = c\_L;}
\NormalTok{outcome\_measures = \{E, A, S, C2, O\_H\} with S = Pr(R \textgreater{}= p) + A {-} C2;}
\NormalTok{C2\_mixed\_search\_minimum\_starts = 3;}
\end{Highlighting}
\end{Shaded}

These controls govern the continuation records of C.2 and C.6 and the
two regression exercises C.6b and C.6c. A price atom is validated by the
rule in \eqref{eq:oa-oa-79} applied to its full preimage, followed by
the four checks of A.11: the buyer's preparation decision is optimal at
its price information set, the conditional target payoff equals the
price pointwise in flow, positive-price and zero-price regions do not
collide (a collision is one preimage), and every unilateral investor
deviation is evaluated against the entire fixed price and preparation
schedule. Average price equal to average payoff is a necessary identity
and is recorded, but it accepts nothing on its own. An exact declared
event is evaluated from its defining identity in fifty-digit arithmetic
with the tie rule applied to the identity; a non-event input has its
sign evaluated at that precision and is reported unresolved if the sign
is closer to zero than the declared resolution. A tolerance for residual
acceptance is not a rule for treating nearby inequalities as equalities.

The outcome measures follow A.7 of the paper appendix. With \(I\) the
preparation indicator and \(V\) the realized challenger value,
\(\mathsf E=\Pr(I=1)\), \(\mathsf A=\Pr(I=1,V\ge p)\),
\(\mathsf S=\Pr(R\ge p\ \text{or}\ [I=1,V\ge p])\), and
\(\mathsf C_2=\Pr(R\ge p,I=1,V\ge p)\). Every continuation row must
satisfy \(\mathsf S=\Pr(R\ge p)+\mathsf A-\mathsf C_2\) by inclusion and
exclusion, with \(\mathsf S\) also integrated directly as a union event,
and \(0\le\mathsf C_2\le\mathsf A\le\mathsf E\le1\),
\(0\le\mathsf S\le1\). High-value ownership \(\mathsf O_H\) is computed
from the actual allocation event; outside the benchmark support it is
not \(e_H/2\) by default, since a high-class value can fail the reserve
or lose to the incumbent.

The flow halfwidth determines a diagnostic mesh, not an integration
truncation. Integrals are evaluated over the full line, or their omitted
tails are explicitly bounded. For a bounded residual \(A\le\overline A\)
and \(|q|\le1\), a Laplace truncation outside \([-T,T]\), \(T>1\), loses
at most \(\overline A e^{-(T-1)/b}\) in per-unit gross profit. Multiply
by the order magnitude for total payoff. Under logistic noise a valid
bound is \(2\overline A/[1+e^{(T-1)/b}]\). Include these bounds in the
error budget; do not compare a tail-truncated integral with a full
integral as though both were exact.

I maintain distinct tolerances for equations, quadrature estimates, and
strategic deviations, and four separate error budgets: residual
approximation, quadrature, tail truncation, and between-grid deviation
coverage. Each is recorded in its own column where the exercise produces
it (\texttt{quadrature\_error}, \texttt{tail\_bound},
\texttt{posterior\_inversion\_error},
\texttt{entry\_independent\_error}, and the deviation-gain estimate with
its rigorous upper bound when one exists); the continuation schema's
\texttt{error\_budget} is the sum of the quadrature and tail terms and
is never a bound on the other two. A solver residual is never reported
as an upper bound on all four. The refinement pass doubles the order
resolution and raises the Gauss-Legendre order on every segment; the
quadrature error estimate on each pass is compared with the declared
target, tightened by a factor of ten on refinement, and a breach is a
breach. If machine precision obstructs the tightened target on some
integral, the check is marked unresolved for that row rather than
passed. Integration warnings are not suppressed. A root or a finite-grid
maximum does not establish equilibrium. For each candidate, evaluate

\[
\begin{aligned}
\epsilon_P&=\sup_{x\ \mathrm{tested}}|P(x)-\mathbb E[V_T\mid X=x]|,\\
\epsilon_e&=\sup_{(P,C)\ \mathrm{tested}}[\text{profit from reversing entry}]_+.
\end{aligned}
\tag{OA.75}\label{eq:oa-oa-75}
\]

and

\[
\epsilon_q=\max_\theta\left\{\sup_{q\in[-1,1]}U_\theta(q;\sigma)-\int U_\theta(q;\sigma)\,d\sigma_\theta(q)\right\}.
\tag{OA.76}\label{eq:oa-oa-76}
\]

The tested price supremum is a numerical diagnostic unless a uniform
bound is supplied. The order supremum must be bounded globally for a
computer-assisted label; a finite approximation is identified
separately. Check all wrong-signed orders as well as correctly signed
ones in numerical diagnostics, even when an analytical sign argument
excludes them. For mixed profiles, every positive-weight support action
must attain the same maximal payoff within the stated tolerance.
Recompute the price and entry schedules once per candidate, then hold
them fixed during each unilateral-deviation calculation.

Mandatory probability identities are \(\int a_\theta=1\),
\(\mathbb E\mu_X=1/2\), and \(\mathbb E P=\mathbb E V_T\). Signal
exercises add the joint-posterior identities in A.7. Independent auction
integration must match closed forms; independent cost-based and
flow-based entry integrations must agree. Repeat numerical diagnostics
after doubling the order resolution and tightening both integration
targets by a factor of \texttt{10}, recorded in the run manifest. A
claimed accepted row that breaches any acceptance bound terminates
validation. Search candidates that fail are retained as rejected rows;
inability to resolve a node is recorded as open, not suppressed or
interpolated.

CSV is the boundary between numerical calculations and presentation.
Files use UTF-8 and quoted headers when a mathematical symbol contains
punctuation. Parameter columns use the symbols
\texttt{h,\ ell,\ p,\ rho,\ c\_L,\ c\_H,\ b,\ k,\ r}; additional columns
such as
\texttt{Delta\_T,\ B\_prior,\ q\_H,\ q\_L,\ e\_H,\ e\_L,\ E,\ O\_H,\ R\_T,\ tau,\ x\_star}
are transliterations of the paper's notation. The data dictionary below
each exercise defines every additional symbol. A branch identifier is a
data label, not an economic selection rule.

The candidate-level error-budget columns shared by C.2 and C.6 are

\begin{Shaded}
\begin{Highlighting}[]
\NormalTok{quadrature\_error, quadrature\_target, quadrature\_error\_refined,}
\NormalTok{quadrature\_target\_refined, quadrature\_refinement, tail\_truncation\_bound,}
\NormalTok{tail\_truncation\_formula, between\_grid\_spacing, between\_grid\_lipschitz\_bound,}
\NormalTok{between\_grid\_closed, between\_grid\_coverage, residual\_approximation\_error,}
\NormalTok{posterior\_inversion\_error, entry\_independent\_error, deviation\_argmax\_H,}
\NormalTok{deviation\_gain\_H, deviation\_argmax\_L, deviation\_gain\_L, error\_budget\_breaches}
\end{Highlighting}
\end{Shaded}

The initial and refined quadrature estimates have separate targets.
\texttt{tail\_truncation\_formula} explains any analytical tail
treatment; zero truncation error is permitted when both infinite tails
are integrated analytically. \texttt{between\_grid\_lipschitz\_bound}
measures the untested deviation allowance, and
\texttt{between\_grid\_closed} states whether that allowance or a
separate argument closes the deviation test.
\texttt{between\_grid\_coverage} names that argument. Residual and
posterior errors remain separate from quadrature error. The two
deviation witnesses identify the best tested order for each investor
type and its gain. \texttt{error\_budget\_breaches} records failed
targets; an unavailable field is \texttt{n/a}.

The quantity registry has columns
\texttt{name,\ value,\ lower,\ upper,\ units,\ display,\ exercise,\ parameter\_set,\ branch,\ status,\ source\_file,\ source\_row,\ definition}.
The \texttt{name} is the exact placeholder key without brackets. Each
name is unique. A missing or failed quantity remains unresolved; it is
never replaced with zero, a cached value from another parameter set, or
a manually typed number. Input declarations can be echoed into the
registry with status \texttt{input}; derived research quantities use the
result-status vocabulary and record their validation. This
administrative input type is not an additional result status.

\subsection{C.1. Baseline, controls, extensions, and scalar
reproduction}\label{oa-c-baseline}

\textbf{Inputs.} Use the benchmark declaration at weak, strong, and
collapse strengths. Evaluate the atomic-cost model with Laplace noise
first. Then replace the noise law by logistic noise, retaining the scale
parameter, and replace cost atoms by a mixture of uniforms on
\([c_L-\varepsilon_C,c_L+\varepsilon_C]\) and
\([c_H-\varepsilon_C,c_H+\varepsilon_C]\). Their mixture weights remain
\(\rho,1-\rho\). These replacements define separate economies.

\textbf{Acquisition layer.} Compute the pointwise auction outcomes from
\eqref{eq:oa-oa-50}. Integrate independently over \(R\), splitting at
\(p\) and \(\ell\), and verify \(t_0,t_H,t_L,g_H,g_L,\Delta_T\). Compute
\(B_r(m),B_r(1/2),B_r(M)\). Theorem margins are

\[
\begin{aligned}
\zeta_L&=B_{r_1}(m)-c_L,\qquad
\zeta_{H0}=c_H-B_{r_0}(1/2),\qquad
\zeta_{H1}=B_{r_1}(M)-c_H,\\
\zeta_0&=k-\Delta_T(r_0),\qquad
\zeta_1=(1-1/b)\rho m\Delta_T(r_1)-k.
\end{aligned}
\tag{OA.77}\label{eq:oa-oa-77}
\]

The collapse node additionally requires \(c_H-B_{r_2}(M)>0\),
\(B_{r_2}(m)-c_L>0\), and its full-order margin. For atomless costs
replace the appropriate cost endpoints as in Proposition A.6. Verify all
margins before assigning an analytical equilibrium label.

\textbf{Financial and entry layer.} Pooling has \(q_H=q_L=0\),
\(\mu=1/2\), \(e=\rho\), and \(P=t_0+\rho[(t_H+t_L)/2-t_0]\). Full
orders use the conditional flow densities and posterior in A.4 or A.6.
Compute \(\tau\), the appropriate threshold, conditional entry,
ownership, mean price, and revenue. Always handle threshold regions
using the cost comparison: a cutoff above the feasible posterior means
no high-cost entry; equality at a Laplace plateau uses the specified tie
rule. Do not insert an infeasible threshold into an interior tail
formula.

For atomless costs, entry at a posterior is the complete mixture CDF
\(H_C(B_r(\mu))\). Independently compute entry by integrating the tail
probability for each cost within each uniform component. The two
integrals must agree. Compute expected paid preparation costs as well as
entry probability; they are not mean cost multiplied by unconditional
entry when the high-cost component is selected.

\textbf{Controls.} The frozen-profile control uses full orders at both
strengths and reoptimizes preparation and prices, but explicitly does
not require that the weak investor chooses that profile. The
price-hidden control reoptimizes investor behavior and pricing, while
the buyer uses the prior: pooling at weak strength and full orders at
strong strength follow their respective bounds. The matched-dividend
control adds exactly
\(\mathcal R_T^{\mathrm{feedback}}-\mathcal R_T^{\mathrm{hidden}}\) to
the hidden financial payoff and price. Verify that all residuals,
orders, and entry probabilities are unchanged. It is not an additional
acquisition payment.

\textbf{Matched-price panel.} The panel distinguishes information from
the mean financial price at the strong strength, for each noise and cost
law in the baseline extension comparison. Its columns are

\begin{Shaded}
\begin{Highlighting}[]
\NormalTok{environment, mean\_financial\_price, seller\_revenue, external\_dividend,}
\NormalTok{preparation\_probability, high\_value\_ownership\_probability, net\_acquisition\_surplus}
\end{Highlighting}
\end{Shaded}

for the three environments \texttt{feedback}, \texttt{price\_hidden},
and \texttt{matched\_dividend} at the strong strength in each of the
four noise and cost-law combinations. The renderer assembles the panel
from the accepted rows of \texttt{tables/equilibrium\_controls.csv} and
\texttt{numerics/feedback\_comparisons.csv}; it solves nothing. The
panel must satisfy, row by row,

\[
D_0=\mathcal R_T^{\mathrm{feedback}}-\mathcal R_T^{\mathrm{hidden}},\qquad
\mathbb E[P^{\mathrm{matched}}]=\mathcal R_T^{\mathrm{hidden}}+D_0=\mathbb E[P^{\mathrm{feedback}}],
\]

The matched row has the same seller revenue, net acquisition surplus,
preparation, high-value ownership, and investor residuals as the
price-hidden environment. The external dividend is zero in the two
equilibrium rows and \(D_0\) in the matched row. Seller revenue in the
matched row does not include the dividend. The dividend is attached to
the traded claim, is not paid by any bidder, and is not a seller
instrument. The panel's note distinguishes the two analytical outcomes
from the matched control, which is an invariance diagnostic and not a
feasible sale mechanism.

\input{tables/table_matched_price.tex}

\textbf{Welfare.} At strong strength, evaluate \eqref{eq:oa-oa-47} or
\eqref{eq:oa-oa-48}, and independently subtract total net allocation
surplus in \eqref{eq:oa-oa-49}. Also compute the revenue difference
directly from conditional entry and from the difference in mean prices.
Store these as different quantities. A strength comparison is not
labeled a welfare effect of access to prices.

\textbf{Validation.} For each equilibrium candidate, integrate residual
profits over the full noise line, splitting at density centers,
posterior kinks, and entry thresholds. Check all orders on the declared
initial and refined grids. Include the special points at which a density
center crosses an entry boundary. The analytical margins provide the
global result for the specified equilibria; the grid is an
implementation check. Verify posterior inversion from the observed
price, including atom probabilities in flat tails. In the frozen-profile
control retain any profitable investor deviations as expected
counterfactual diagnostics rather than mistakenly treating that profile
as equilibrium.

\textbf{Outputs.}

\begin{Shaded}
\begin{Highlighting}[]
\NormalTok{tables/auction\_primitives.csv:}
\NormalTok{parameter\_set, r, p, t\_0, t\_H, t\_L, g\_H, g\_L, Delta\_T, B\_m, B\_prior, B\_M}

\NormalTok{tables/equilibrium\_controls.csv:}
\NormalTok{parameter\_set, noise, cost\_law, r, experiment, q\_H, q\_L,}
\NormalTok{e\_H, e\_L, E, O\_H, R\_T, W, expected\_preparation\_cost,}
\NormalTok{tau, x\_star, matched\_dividend, status, accepted}

\NormalTok{tables/extensions.csv:}
\NormalTok{parameter\_set, noise, cost\_law, r\_weak, r\_strong,}
\NormalTok{E\_weak, E\_strong, O\_H\_weak, O\_H\_strong,}
\NormalTok{zeta\_L, zeta\_H0, zeta\_H1, zeta\_0, zeta\_1, status, accepted}

\NormalTok{numerics/baseline\_deviations.csv:}
\NormalTok{parameter\_set, experiment, noise, cost\_law, r, state, q,}
\NormalTok{U(q), U(candidate), deviation\_gain, quadrature\_error, tail\_bound}

\NormalTok{numerics/feedback\_comparisons.csv:}
\NormalTok{parameter\_set, noise, cost\_law, r,}
\NormalTok{W\_feedback, W\_hidden, W\_gain,}
\NormalTok{R\_T\_feedback, R\_T\_hidden, R\_T\_gain, matched\_dividend,}
\NormalTok{welfare\_identity\_error, revenue\_identity\_error,}
\NormalTok{residual\_invariance\_error, status, accepted}

\NormalTok{figures\_data/two\_returns.csv:}
\NormalTok{r, mu, Delta\_T, B\_r(mu), d\_Delta\_T\_dr, d\_B\_r\_mu\_dr}

\NormalTok{tables/matched\_price.csv (derived by the table renderer):}
\NormalTok{parameter\_set, noise, cost\_law, r, environment, mean\_financial\_price,}
\NormalTok{seller\_revenue, external\_dividend, preparation\_probability,}
\NormalTok{high\_value\_ownership\_probability, net\_acquisition\_surplus, role, status, accepted}
\end{Highlighting}
\end{Shaded}

\(W\) denotes allocation value net of paid preparation costs, not target
revenue. The derivatives in the last file are the explicit uniform
formulas. Use the correspondence strength grid from C.2 and beliefs
\(m,1/2,M\). These outputs feed Tables 1--3 and Figure 1. Scalar keys
are defined in C.8; no displayed value is copied directly from a figure.

The complete extension margins accompany the preparation comparisons in
the main paper. A negative sufficient-condition margin does not reject a
separately validated outcome.

\input{tables/table_extensions_margins.tex}

\subsection{C.2. The equilibrium correspondence and interval
certificates}\label{oa-c-correspondence}

\textbf{Inputs and grid.} Use the benchmark primitives. The initial
strength mesh is the exact decimal arithmetic progression from
\texttt{1.005} through \texttt{3.800} by \texttt{0.005}. Add the weak,
strong, collapse, and certified strengths and the boundaries
\(\mathfrak r(k),r_N,r_U,r_C\), with offsets \texttt{0.0001} and
\texttt{0.001} on both sides. Boundary formulas are evaluated
numerically to select nodes; the declared equality at \(r_C\) is
evaluated symbolically as \(\tau=M\) for preparation on the upper
plateau. Adjacent strength intervals whose accepted branch sets differ
are refined to spacing \texttt{0.001}. The run manifest records the
nodes and refinement intervals.

The certificate input brackets are exact decimals:

\begin{Shaded}
\begin{Highlighting}[]
\NormalTok{r = 1.55; v\_left = 0.46031618; v\_right = 0.46031620;}
\NormalTok{r = 1.60; v\_left = 0.70747537; v\_right = 0.70747539;}
\NormalTok{r = 1.65; v\_left = 0.90333198; v\_right = 0.90333201.}
\end{Highlighting}
\end{Shaded}

\textbf{Analytically classified outcomes.} First verify membership in
the maintained floor/prior domain at each node. Pooling existence is
given by \(k-\rho\Delta_T/2\ge0\) there; its sufficient uniqueness bound
is \(k-\Delta_T>0\). Full-order uniqueness is established by
\((1-1/b)\rho m\Delta_T-k>0\). Record these margins individually.
Compute \(r_N,r_U,r_C\) from A.5 and verify each defining equation,
including support restrictions. Outside a maintained domain, evaluate
the actual prior entry rule and continuation instead of retaining a
classification based on \(\rho\).

\textbf{Full-order candidates.} At each node, form the full-order
posterior and preparation policy. The uniform derivative bound can
establish uniqueness, and the sharper candidate-specific \(J\) test in
\eqref{eq:oa-oa-22} can establish existence. A failed sufficient test
does not reject the profile. Evaluate both traders on the initial and
refined deviation grids, including wrong-signed orders, and record any
profitable deviation with its action. For the pure low-type residual,
strict concavity reduces boundary optimality to \(U_L'(1)\ge0\). An
unenclosed high-type derivative mesh is a diagnostic. Acceptance without
an analytical argument or interval cover remains a numerical diagnostic,
even if every tested deviation is unprofitable.

\textbf{Asymmetric continuation.} At each \(r\), solve \(\Psi(r,v)=0\)
for candidates with \((1,-v)\) over the searched magnitude interval
\texttt{{[}0.005,\ 0.9999{]}}. Bracket sign changes on a mesh of spacing
\texttt{0.01}, add the declared certified brackets, and solve bracketed
roots by Brent's method. Minimize \(|\Psi|\) near observed local minima
and subdivide each such neighborhood on a 41-point mesh to seek
additional sign changes. A sign change at a discontinuity without a
vanishing residual does not establish a root. A validated tangency
candidate remains a numerical diagnostic; a failed or unresolved attempt
retains its witness. The unilateral derivative in \(s\) holds the
candidate schedule fixed, whereas changing \(v\) changes that schedule.

At each candidate compute \(x^*,e_H,e_L,\mathsf E,\mathsf O_H\),
revenue, the pooling-existence flag, and the deviation diagnostics. At
the three declared nodes, \hyperref[oa-b]{Appendix B} encloses the root,
proves its endpoint signs, verifies low-type concavity, and covers
high-type deviations uniformly over the root bracket. No other node
inherits those certificates. An interior root approaching \(v=1\) is a
branch boundary; the boundary profile is evaluated separately.

\textbf{Other pure and mixed candidates.} The complete correspondence
remains open. The pure search evaluates \((u,-v)\) on \([0,1]^2\) from
the nine initial pairs in \(\{0.2,0.6,0.95\}^2\), with forty damped
best-response iterations and a local fixed-point solve. Interior pure
roots are polished using both traders' held-schedule first-order
conditions before revalidation. Every resulting candidate faces the full
initial and refined deviation scans. The mixed search is restricted to
high-type mixtures over correctly signed orders and a pure low-type
magnitude \(v\in[0,1]\). It runs meshes of spacing \texttt{0.05} and
\texttt{0.025}. This restriction defines the search domain; it is not an
exclusion of low-type mixing in every equilibrium. A mesh check of
residual monotonicity is recorded, but does not supply the global
premise of the Stieltjes lemma for an arbitrary mixed profile.

A mixed candidate with support points \(q_{\theta j}\) and weights
\(\omega_{\theta j}\) has conditional density
\(a_\theta(x)=\sum_j\omega_{\theta j}f(x-q_{\theta j})\). Recompute its
posterior, entry, and competitive price. Its payoff is the weighted
average of the support payoffs. All positive-weight points must be best
responses within tolerance, and every off-support deviation must be
checked. Preserve distinct roots and supports found from different
initializations. A candidate search that fails to resolve a node
receives an open flag; it does not create an empty equilibrium set.

\textbf{Mixed-search attempt records.} On each mesh, three starts place
high-type weights uniformly, near the full purchase, and near zero, with
initial low-type magnitudes \texttt{0.5}, \texttt{0.9}, and
\texttt{0.1}. Damped multiplicative-weight updates stop on convergence,
a detected stall, or the declared iteration budget, whose default is
\texttt{120}. The final support is then subjected to payoff-indifference
and simplex equations, with bounded support additions and removals. A
collapsed pure profile is polished by the continuous pure solver and
revalidated before deduplication. Each start appears in
\texttt{numerics/mixed\_search\_attempts.csv}, including its initial
values, domain, budget, stopping reason, removed support points and
their residuals, and final support.
\texttt{numerics/mixed\_supports.csv} records each state's final support
action, weight, payoff, and gap to the best tested deviation. Any
unresolved start gives the node an explicit open search row in
\texttt{numerics/correspondence.csv}; that row does not invalidate a
separately established equilibrium there. None of the recorded search
outcomes is a nonexistence theorem.

\textbf{Continuation identity and deduplication.} Two candidates at a
node are the same continuation only when their complete strategies
(support points and weights within the order identity tolerance) and
their induced price information (every positive-mass price atom's value,
state masses, and kind within the atom identity tolerance) agree,
together with the preparation rule; this is the rule of
\eqref{eq:oa-oa-53} applied numerically, and C.6 uses the same one. On
the benchmark strength grid the low-cost floor \(c_L<B_r(m)\) holds at
every node, entry is positive at every flow, and no price pool combines
distinct raw posteriors. Constant-posterior atoms remain possible, and
their prices are fixed by the orders, so identity reduces to the
complete order profile there; the code applies the full rule regardless.
Strength labels are normalized to their shortest exact decimal before
nodes are compared, so a strength written two ways is one node. A root
found by \(|\Psi|\) minimization is compared with the roots found by
sign change at the same node and merged when it is the same
continuation; the raw attempt stays in the ledger with the identifier of
the row that absorbed it. \texttt{multiplicity\_found} counts distinct
accepted continuations after this merge, never raw candidates. Rejected
candidates keep their breach witness; a candidate rejected because a
sufficient condition failed is not a rejection, and the status
vocabulary of C.0 applies.

\textbf{Figure discipline.} Plot analytical uniqueness regions
separately from existence-only regions. Show every accepted observed
branch. Certified points have interval bars. Numerical lines are broken
at failed validation, discontinuities, branch changes, or unresolved
gaps. The multiplicity flag means that distinct equilibria have been
established or found at a node, not that all equilibria have been
enumerated. The ordered certified entry intervals provide the stated
cross-economy comparison; a fitted derivative on an exploratory line is
not substituted for that result.

\textbf{Outputs.}

\begin{Shaded}
\begin{Highlighting}[]
\NormalTok{numerics/correspondence.csv and numerics/correspondence\_attempts.csv:}
\NormalTok{r, branch, q\_H, q\_L, v, e\_H, e\_L, E, O\_H, R\_T, tau, x\_star, pooling\_exists,}
\NormalTok{pooling\_unique\_bound, full\_unique\_bound, existence\_status, uniqueness\_status, accepted,}
\NormalTok{multiplicity\_found, epsilon\_P, epsilon\_e, epsilon\_q, tail\_bound, unresolved\_reason,}
\NormalTok{candidate\_id, continuation\_id, parameter\_set\_id, search\_path, duplicate\_of,}
\NormalTok{n\_distinct\_accepted, result\_status}
\NormalTok{plus the shared error{-}budget columns in C.0}

\NormalTok{numerics/mixed\_supports.csv:}
\NormalTok{r, branch, state, support\_index, q, weight, U(q), support\_gap, off\_support\_gain\_mesh,}
\NormalTok{accepted, status, init\_id, mesh\_spacing, iterations, stop\_reason, final\_gap,}
\NormalTok{gap\_to\_best\_tested, simplex\_residual, residual\_monotone, quadrature\_error,}
\NormalTok{search\_domain, outcome, candidate\_id, continuation\_id, parameter\_set\_id}

\NormalTok{numerics/mixed\_search\_attempts.csv:}
\NormalTok{r, mesh, start\_id, start\_description, v\_init, domain, budget, iterations, stop\_reason,}
\NormalTok{final\_gap, final\_v\_residual, support\_solve, support\_solve\_residual, support\_rounds,}
\NormalTok{pruned\_points, support, support\_payoffs, gap\_to\_best\_tested, best\_tested\_q, v,}
\NormalTok{simplex\_residual, min\_weight, residual\_monotone, outcome, witness, low\_support\_payoff,}
\NormalTok{low\_gap\_to\_best\_tested, validation}

\NormalTok{numerics/certificates.csv:}
\NormalTok{r, v\_lower, v\_upper, Psi\_left\_lower, Psi\_left\_upper,}
\NormalTok{Psi\_right\_lower, Psi\_right\_upper, Gamma\_H\_lower,}
\NormalTok{L\_U\_upper, mesh\_intervals, interval\_digits,}
\NormalTok{E\_lower, E\_upper, pooling\_margin\_lower, threshold\_margin\_lower, accepted}

\NormalTok{numerics/thresholds.csv:}
\NormalTok{boundary, value, lower, upper, defining\_residual,}
\NormalTok{in\_support\_domain, low\_cost\_floor\_valid, interpretation}
\end{Highlighting}
\end{Shaded}

\texttt{correspondence\_attempts.csv} keeps raw candidate attempts and
their \texttt{duplicate\_of} links. \texttt{correspondence.csv} removes
duplicate attempts while preserving rejected and open records.
\texttt{search\_path} describes the candidate construction;
\texttt{n\_distinct\_accepted} and \texttt{multiplicity\_found} count
distinct accepted continuations at the node. The support file links each
validated strategy to the full candidate and continuation identity;
unvalidated starts carry \texttt{n/a} for those foreign keys. The
attempt ledger identifies \texttt{converged-candidate},
\texttt{converged-pure}, and \texttt{unresolved} search outcomes
separately from validation status. \texttt{support\_payoffs},
\texttt{gap\_to\_best\_tested}, \texttt{pruned\_points}, and
\texttt{witness} preserve the reasons for retaining, removing, or
rejecting a support.

The threshold file uses exact row labels
\texttt{pooling\_unique\_sufficient}, \texttt{pooling\_existence},
\texttt{full\_orders\_unique\_sufficient}, and
\texttt{high\_cost\_ceiling}. It also carries separately labeled
analytical scalars \texttt{m}, \texttt{M}, and
\texttt{laplace\_entry\_left\_limit} needed by the registry; these are
not mislabeled activation thresholds. Every row stores its defining
expression and checks. Fields that do not apply to a probability scalar
are marked not applicable, not set to zero.

All general parameter columns are included or keyed to an immutable
parameter declaration whose full contents accompany the file.
\texttt{uniqueness\_status} distinguishes analytical uniqueness from not
established; it is never inferred from the number of search hits.
\texttt{off\_support\_gain\_mesh} is a mesh maximum and is named as one;
it is not a bound. \texttt{quadrature\_error} is checked against the
declared target on each pass, and \texttt{tail\_bound} is zero under
Laplace noise only because the exterior tails are integrated in closed
form against the constant residual. Feed: Figure 2, Proposition 3, the
threshold discussion, and the baseline table. Registry keys
\texttt{cert\_a\_*}, \texttt{cert\_b\_*}, and \texttt{cert\_c\_*} refer
to the ordered declared nodes, not arbitrary roots selected from a
larger search.

\subsection{C.3. Complementary private signals}\label{oa-c-signals}

\textbf{Inputs.} Start with the complementary-signal declaration. Then
cross trader accuracies \texttt{\{0.68,\ 0.69,\ 0.70,\ 0.71,\ 0.72\}}
with buyer accuracies \texttt{\{0.73,\ 0.74,\ 0.75,\ 0.76,\ 0.77\}},
holding other primitives fixed. These are separate equilibria with exact
decimal parameters, not a claim that every grid point meets Proposition
A.7.

\textbf{Objects.} Compute \(m,M\), public fundamental bounds
\(\mu_-,\mu_+\), the worst joint buyer posterior \(\phi_-(\mu_-)\), and
the best joint posterior \(\phi_+(\mu_+)\). Compute the five strict
margins from \eqref{eq:oa-oa-29}: the low-cost floor, high-cost
exclusion using the buyer's favorable private signal at weak strength,
high-cost profitability after a favorable strong public price and
private signal, weak residual exclusion, and strong global derivative
margin. If a margin fails, withhold the analytical theorem label; do not
conclude that the economic reversal is absent.

At both strengths construct signal-contingent order laws and the correct
state-dependent entry policy. For the theorem-supported weak node public
belief is the prior and both high-cost private-signal types stay out.
For the strong node use full orders by \(T\), then compute the public
posterior from \(\lambda_X\), each buyer posterior \(\phi_Y(\mu_X)\),
and both private-signal cutoffs in \eqref{eq:oa-oa-39}. Always handle
both \(Y=+\) and \(Y=-\), including always/never and plateau-equality
cases. The benchmark example's exclusion of one private-signal type is
an outcome, not an imposed simplification for the sweep.

\textbf{Independent verification.} First integrate entry using
\eqref{eq:oa-oa-40}--\eqref{eq:oa-oa-41}. Independently sum over the
finite fundamental, trader-signal, and buyer-signal states and integrate
\(Z\) and \(C\). Compute the price directly as the conditional
expectation of \(V_T\) given flow and compare it with
\eqref{eq:oa-oa-35}. Recover the public posterior from the price
inversion and then update using \(Y\); compare with the direct joint
likelihood. Verify the residual identities \eqref{eq:oa-oa-37} by direct
conditional trader payoffs. Finally evaluate all unilateral orders using
the fixed candidate schedules. The investor must not be given the
buyer's private signal during this calculation.

\textbf{Outputs.}

\begin{Shaded}
\begin{Highlighting}[]
\NormalTok{numerics/two\_signals.csv:}
\NormalTok{a, d, r, q\_plus, q\_minus, mu\_lower, mu\_upper,}
\NormalTok{phi\_minus\_mu\_lower, phi\_plus\_mu\_upper,}
\NormalTok{x\_star\_Yplus, x\_star\_Yminus, e\_H, e\_L, E, O\_H, R\_T,}
\NormalTok{low\_cost\_margin, private\_only\_exclusion\_margin, joint\_entry\_margin,}
\NormalTok{weak\_order\_margin, strong\_order\_margin,}
\NormalTok{posterior\_error, residual\_error, epsilon\_P, epsilon\_q,}
\NormalTok{status, accepted}

\NormalTok{numerics/two\_signal\_deviations.csv:}
\NormalTok{a, d, r, trader\_signal, q, U(q), U(candidate),}
\NormalTok{deviation\_gain, quadrature\_error, tail\_bound}
\end{Highlighting}
\end{Shaded}

The same complete primitive vector accompanies each row. The
scalar-margin names map to these columns in order: low cost to
\texttt{low\_cost\_margin}, high prior to
\texttt{private\_only\_exclusion\_margin}, high ceiling to
\texttt{joint\_entry\_margin}, weak trade to
\texttt{weak\_order\_margin}, and strong trade to
\texttt{strong\_order\_margin}. These are comparison-level margins;
repeat them consistently on both strength rows or store them in a
separately keyed comparison record. The display accuracy keys are
percentages generated from the raw probability keys, never separately
typed. Table 3 in the main paper uses the declared example. The signal
table below reports every accuracy combination and distinguishes the
sufficient uniqueness region from other validated outcomes. These
outputs also supply the example of Proposition A.7.

The full signal comparison is reported below. The declared example is
marked; every other accuracy combination remains visible.

\input{tables/table_signal_grid.tex}

\subsection{C.4. Moderate values and constructive
nonemptiness}\label{oa-c-moderate}

\textbf{Inputs.} Reproduce the moderate declaration exactly. Compute the
same primitives, theorem margins, complete equilibrium objects, and
controls as in C.1. Use independent auction integration and global
analytical margins, not only a best-response mesh, to validate the
theorem-supported example.

For the nonemptiness construction, set \texttt{ell\ =\ 1},
\texttt{p\ =\ 0.5}, \texttt{rho\ =\ 0.25}, \texttt{b\ =\ 2}, and use
high-to-low ratios \texttt{\{1.05,\ 1.10,\ 1.25,\ 1.50,\ 2.00\}}. For
each ratio let \(D=h-\ell\) and examine \(\delta_j=D2^{-j}\) for integer
\(j\) from \texttt{1} through \texttt{20}. Define \(r_1=\ell+\delta_j\)
and \(r_0=\ell+\delta_j^2/D\). Test the conditions needed in A.6. When
the intervals are nonempty choose

\[
\begin{aligned}
k&=\frac{\Delta_T(r_0)+(1-1/b)\rho m\Delta_T(r_1)}2,\\
c_H&=\frac{B_{r_0}(1/2)+B_{r_1}(M)}2,
\qquad c_L=\frac{B_{r_1}(m)}2.
\end{aligned}
\tag{OA.78}\label{eq:oa-oa-78}
\]

Every chosen value is an output of the construction, not a fixed
calibration imposed on all ratios. Check \(c_L<c_H\), support, and all
theorem inequalities explicitly. Keep the sequence of failed and
successful \(j\) values. Small margins require adequate arithmetic
precision; floating-point cancellation is not evidence against
nonemptiness. Report the first certified feasible member of the declared
sequence and the numerical scale of \(k\), not an invented uniform lower
bound on its size.

\textbf{Outputs.}

\begin{Shaded}
\begin{Highlighting}[]
\NormalTok{numerics/moderate\_values.csv:}
\NormalTok{h, ell, p, rho, c\_L, c\_H, b, k, r\_weak, r\_strong,}
\NormalTok{zeta\_L, zeta\_H0, zeta\_H1, zeta\_0, zeta\_1,}
\NormalTok{E\_weak, E\_strong, O\_H\_weak, O\_H\_strong, status, accepted}

\NormalTok{numerics/nonemptiness\_construction.csv:}
\NormalTok{h\_over\_ell, j, delta, r\_weak, r\_strong, k, c\_L, c\_H,}
\NormalTok{all\_margins\_min, feasible, arithmetic\_precision, reason}
\end{Highlighting}
\end{Shaded}

Feed: Table 3 and the moderate-value paragraph. The general construction
is an analytical argument; the finite sequence is a diagnostic
illustration of its scale.

\subsection{C.5. Noise laws, threshold distance, and posterior upper
tails}\label{oa-c-noise}

\textbf{Inputs.} Fix benchmark strong strength, \(b\), and \(\rho\). Use
full orders. Form \texttt{401} equally spaced points for the normalized
threshold distance \(d_\tau=(M-\tau)/(M-1/2)\) on \([0,1]\), and add the
benchmark \(\tau\), \(M-10^{-4}\), \(M-10^{-6}\), and \(M+10^{-6}\). At
the endpoints use the stated mathematical limits and tie convention
rather than unstable logarithms. Compare Laplace and logistic laws at
the same scale parameter; do not claim equal variance or a Blackwell
ordering between the laws.

\textbf{Objects.} For each threshold, compute the order-flow cutoff,
\(\alpha_H(\tau),\alpha_L(\tau)\), posterior upper-tail mass
\([\alpha_H+\alpha_L]/2\), and the induced entry
\(\rho+(1-\rho)[\alpha_H+\alpha_L]/2\). Verify each tail by direct
integration of the conditional flow densities. For the baseline logistic
example, compute \(x^*/(b\pi/\sqrt3)\) from the derived noise variance,
not from the aggregate-flow variance. Store both variances separately if
both standardized cutoffs are reported.

A threshold scan is first an information-experiment comparison. To
interpret a row as an equilibrium with changed preparation cost, set
\(c_H(\tau)=g_L+\tau(g_H-g_L)\), keep the other primitives fixed, and
recheck the low-cost floor, cost ordering, and global full-order bound.
Label only validated rows as such. At the strong benchmark the
full-order bound does not depend on the high-cost threshold, but that
fact must not be presumed after changing other primitives.

At \(\tau=M\), Laplace admits entry in its entire upper posterior
plateau under the tie convention; logistic has zero mass at the
unattained endpoint. For \(\tau>M\), both tail masses are zero. This
distinction is part of the plotted result. On the lower endpoint
\(\tau=1/2\), the threshold is zero and no artificial singularity
arises.

\textbf{Outputs.}

\begin{Shaded}
\begin{Highlighting}[]
\NormalTok{figures\_data/posterior\_tails.csv:}
\NormalTok{noise, b, tau, M\_minus\_tau, normalized\_distance, x\_star,}
\NormalTok{alpha\_H, alpha\_L, posterior\_upper\_tail\_mass, E,}
\NormalTok{implied\_c\_H, noise\_variance, flow\_variance, threshold\_noise\_sd,}
\NormalTok{full\_order\_margin, equilibrium\_interpretation\_valid,}
\NormalTok{integration\_error, status}
\end{Highlighting}
\end{Shaded}

Feed: Figure 3, the logistic magnitude paragraph, and Table 3's noise
rows. Infinite thresholds are stored with an explicit
\texttt{unattainable} label, not converted to an arbitrary finite
cutoff.

\subsection{C.6. Reserve comparisons and exploratory seller
continuation}\label{oa-c-reserve}

\textbf{Inputs.} Reproduce the binary-value comparison using the
benchmark reserves \texttt{0.5} and \texttt{1.01} at strengths
\texttt{1.2} and \texttt{3}. Reproduce the atomless-value class economy
using half-width \texttt{0.05} and reserves \texttt{0.5} and
\texttt{1.1} at the same strengths. In the class economy the investor
sees the class only and preparation reveals the exact value. Do not let
the investor trade on the within-class realization.

\textbf{Payoff oracle.} Integrate the realized outcomes
\eqref{eq:oa-oa-50} over the independent uniform incumbent and each
conditional value distribution. Split at \(R=p\), \(R=v\), and all
reserve/value support boundaries. Verify against \eqref{eq:oa-oa-51}
and, where applicable, \eqref{eq:oa-oa-52}. Compute class-conditional
\(t_H,t_L,g_H,g_L\) for every reserve, including reserves inside a value
band. A formula for \(p<\ell\) or for exclusion of an entire class is
not used inside a partially excluded band. Each reserve's regime on the
full domain (below \(\ell\), at \(\ell\), between \(\ell\) and \(r\), at
\(r\), between \(r\) and \(h\), at \(h\), above \(h\), and the band
cases for the class economy) is decided by exact comparison of the
declared numbers and written to the row.

\textbf{Complete continuations.} A candidate at a reserve is accepted
only as a complete continuation in the sense of \eqref{eq:oa-oa-53}: the
investor's state- or class-contingent order distribution, the pricing
rule with every positive-mass price pool, the buyer's information at
each price (the pooled posterior on an atom, the inverted posterior
elsewhere), the preparation rule by observed price and cost, and the
validation evidence. For a candidate order profile, compute its raw flow
posterior \(\mu_X\). When \(\Delta_T>0\) and there is positive entry,
the benchmark inversion applies to the positive-entry prices. If entry
vanishes on a set of flows, those flows pool at \(P=t_0\), and their
posterior is the integral over the full preimage of that price:

\[
\mu(P=t_0)=\frac{\int_{\{x:P(x)=t_0\}}a_H(x)\,dx}
{\int_{\{x:P(x)=t_0\}}[a_H(x)+a_L(x)]\,dx}.
\tag{OA.79}\label{eq:oa-oa-79}
\]

The formula is used only when the denominator is positive. The buyer's
entry rule at that price uses this pooled posterior and nothing finer; a
candidate whose buyer acts on the raw posterior inside the pool is not
measurable with respect to the price and is rejected. The convenient
construction for independent costs,
\(P(\mu)=t_0+H_C(B(\mu))[t_L-t_0+\Delta_T\mu]\), prices every
positive-entry flow by the inversion and pools exactly the zero-entry
preimage; it is strictly increasing wherever entry is positive and
\(\Delta_T>0\). It is the one pricing family the C.6 node solver
constructs, and the row says so: \texttt{pricing\_family\_coverage}
records that alternative cutoffs were not scanned at that node and
\texttt{exhaustive\_pricing\_search} is false.
\hyperref[oa-a-pools]{Section A.11} shows why this matters. In the
example of A.11, the same orders admit a family of pools, each a
different continuation, and the convenient construction is one member.
This example does not establish such a family at every reserve. Every
accepted candidate must pass the price-atom validation of C.0: pooled
posterior from the full preimage, buyer optimality at the pooled
posterior and pointwise on positive prices, competitive pricing
pointwise in flow on both regions, no collision between the atom and the
positive-price range, and every unilateral deviation against the whole
fixed schedule.

At reserves with \(\Delta_T=0\), all entry impossible, or a degenerate
payoff, analyze the constant-price candidate directly. Pooling uses the
actual prior entry probability \(e_0=H_C(B(1/2))\), and its correctly
signed deviation coefficient is \(e_0\Delta_T/2-k\); it is not
automatically \(\rho\Delta_T/2-k\). When \(e_0=0\), no trade and no
entry are supported by constant pricing. All accepted candidates must
satisfy the original conditional pricing and entry conditions.

\textbf{Specified comparisons.} At the declared reserves, verify the
relevant strict no-trade or full-order bounds using the class payoff
spread and the actual floor \(e(m)=H_C(B_r(m))\). Compute entry and
seller revenue and compare each alternative reserve with the original
one at the same strength. This reproduces the profitable alternatives,
not a global reserve optimum.

\input{tables/table_reserve_details.tex}

\textbf{Event grid.} The sweep and the event exercise share one grid
rule. Beyond the declared reserve mesh, every node set contains the
support and participation events \(0\), \(\ell\), \(r\), \(h\); in the
class economy \(\ell\pm\varepsilon_V\) and \(h\pm\varepsilon_V\); the
declared reserve alternatives; and the two exact events of the paper
appendix, the low-cost floor equality \(p_L=h-c_L/m\) and the
expensive-entry ceiling equality \(p_H=\sqrt{2r(h-c_H/M)-r^2}\), each
attached to the strength at which its defining regime applies. Every
event carries one-sided offsets \texttt{1e-4}, \texttt{1e-6}, and
\texttt{1e-8} on both sides where the offset stays in \([0,h]\) (or
\([0,h+\varepsilon_V]\)), so that the regime classification on each side
is tested separately. An exact event is stored by its defining relation
and recomputed from it at fifty digits; its decimal expansion is
reference only and never decides a tie. The tie rule is applied to the
identity: at \(p_L\) the low-cost type enters on the lower plateau under
\(B(m)=c_L\), and at \(p_H\) the high-cost type enters on the upper
plateau under \(\tau=M\) with \(x^*=1\). A node whose inequality sign
cannot be resolved at that precision is written as unresolved. The same
continuation need not be supported on both sides of an event, and the
record does not assume it.

\textbf{Exploratory sweep.} The initial binary-value reserve mesh runs
from \texttt{0} to \texttt{h} in steps of \texttt{0.05}; the class-value
mesh runs to \texttt{h\ +\ epsilon\_V}, with the event nodes added in
both cases. Intervals with changing accepted branch sets, an unresolved
endpoint, or a found revenue maximum are refined to spacing
\texttt{0.002}. At each applicable node the solver checks pooling, full
orders, asymmetric roots and tangencies, and pure profiles from the nine
starts in \(\{0.2,0.6,0.95\}^2\), with forty damped best-response
iterations and a local fixed-point solve. Its mixed search uses a
\texttt{0.05} high-type order mesh, initially uniform weights, an
initial pure low-type magnitude \texttt{0.5}, and at most \texttt{150}
multiplicative-update iterations with tolerance \texttt{1e-7}. This
single restricted mixed search differs from the multiple-start C.2
exercise and supplies no claim about unsearched supports. Candidate
schedules use the pricing family described above. Initializations and
outcomes are retained in the attempt ledger; there are no bidirectional
continuation warm starts. Analytical uniqueness can make the remaining
candidate searches unnecessary at a node, and the manifest records that
reason.

\textbf{Identity and deduplication.} Every candidate is written to the
ledger with a \texttt{candidate\_id} that records how it was obtained
(exercise, node, branch, and sequence) and a \texttt{continuation\_id}
that is the canonical economic identity: the parameter-set identity
built from the exact declarations, the institution, the information
structure, the tie rule, the complete order profile at the order
identity tolerance, the pricing family, the preparation rule, and every
positive-mass price atom at the atom identity tolerance. Atoms and pools
below the numerical probability cutoff are omitted from the numerical
identity comparison. This computational cutoff does not make a
positive-probability event a mathematical null set; its mass remains
subject to the probability error budget. Accepted candidates are merged
only when their continuation identities agree; the absorbed candidate
keeps its row with \texttt{duplicate\_of} pointing at the
representative. Candidates are never merged on a rounded strength or
reserve, on orders or mixed supports alone, on entry probability alone,
on mean price or revenue alone, or on a solver's branch name. Two
candidates with the same orders and different price atoms stay distinct,
and the row's \texttt{identity\_note} says so. Counts of continuations
found, rejected candidates, unresolved candidates, and merged duplicates
are computed from this ledger, never from the grid length.

\textbf{Reporting.} For each reserve retain every accepted continuation
found. Report the minimum and maximum entry and revenue across those
continuations and the search outcome, which is one of: analytically
unique continuation; accepted continuations found with uniqueness not
established; no accepted candidate found by the declared searches, which
is not a nonexistence proof; or unresolved search with the open rows
retained. These are found-continuation ranges, not certified envelopes
of all equilibria. Never replace a failed solve with the nearest
successful reserve, optimize a frozen information experiment, or label
the best sampled reserve globally optimal. Any proposed seller
equilibrium must specify its continuation at unchosen reserves as well
as chosen ones.

The online summary of the sweep is headed ``highest revenue among
continuations found'', never ``optimal reserve''. It lists, per economy
and strength, the reserve and continuation with the highest accepted
revenue, the number of reserve nodes attempted, the number with more
than one distinct accepted continuation, and the number of unresolved
nodes, all computed from the attempt ledger. It carries its known limits
in the same table note: the search covers the declared candidate classes
only, the pricing family at each node is the benchmark construction, and
no envelope is certified. It is not in the paper.

The coverage counts use the final pass at each reserve. Earlier attempts
that triggered continuation-identity refinement remain in the numerical
archive, with their inputs and outcomes; they are not counted as
additional final-pass nodes. A bounded refinement of near-root
candidates is followed by full revalidation before economically
identical continuations are merged.

\input{tables/table_reserve_exploratory.tex}

\textbf{Outputs.}

\begin{Shaded}
\begin{Highlighting}[]
\NormalTok{tables/reserve\_comparisons.csv:}
\NormalTok{value\_law, signal\_information, epsilon\_V, r, p,}
\NormalTok{t\_0, t\_H, t\_L, g\_H, g\_L, Delta\_T,}
\NormalTok{q\_H, q\_L, e\_H, e\_L, E, R\_T,}
\NormalTok{low\_cost\_floor\_margin, trading\_margin, payoff\_oracle\_error,}
\NormalTok{status, accepted}

\NormalTok{numerics/reserve\_continuations.csv (and numerics/reserve\_events.csv):}
\NormalTok{value\_law, r, p, branch, q\_H, q\_L,}
\NormalTok{E, O\_H, R\_T, no\_entry\_price\_mass, posterior\_in\_no\_entry\_pool,}
\NormalTok{epsilon\_P, epsilon\_e, epsilon\_q, status, accepted, unresolved\_reason,}
\NormalTok{p\_exact, regime, band\_status, low\_cost\_floor\_relation, ceiling\_relation,}
\NormalTok{duplicate\_of, identity\_note,}
\NormalTok{candidate\_id, continuation\_id, parameter\_set\_id, institution\_id,}
\NormalTok{information\_structure\_id, order\_rule\_id, price\_rule\_id,}
\NormalTok{pricing\_family\_coverage, exhaustive\_pricing\_search,}
\NormalTok{pool\_id, pool\_set\_definition, pool\_cutoff\_exact, price\_atom\_value,}
\NormalTok{state\_H\_pool\_mass, state\_L\_pool\_mass, pool\_probability, pool\_posterior,}
\NormalTok{preparation\_rule\_id, tie\_rule\_id, event\_id, event\_defining\_relation,}
\NormalTok{preparation\_probability, admissible\_challenger\_probability, sale\_probability,}
\NormalTok{two\_admissible\_bidders\_probability, high\_value\_ownership\_probability,}
\NormalTok{seller\_revenue, mean\_financial\_price,}
\NormalTok{pricing\_error, belief\_error, entry\_deviation\_gain,}
\NormalTok{investor\_deviation\_gain\_estimate, investor\_deviation\_gain\_upper\_bound,}
\NormalTok{error\_budget, result\_status, existence\_scope, uniqueness\_scope,}
\NormalTok{search\_coverage\_scope, rejection\_reason, run\_id}
\NormalTok{plus the shared error{-}budget columns in C.0}

\NormalTok{numerics/reserve\_ranges.csv (and numerics/reserve\_event\_ranges.csv):}
\NormalTok{value\_law, r, p, accepted\_continuations\_found,}
\NormalTok{E\_min\_found, E\_max\_found, R\_T\_min\_found, R\_T\_max\_found,}
\NormalTok{search\_unresolved, global\_envelope\_certified,}
\NormalTok{p\_exact, event\_id, candidates\_evaluated, candidates\_accepted\_raw,}
\NormalTok{candidates\_rejected, candidates\_unresolved, duplicates\_merged, search\_outcome}
\end{Highlighting}
\end{Shaded}

\textbf{Data dictionary for the continuation schema.} \texttt{p} is the
exact declared decimal, or the thirty-two-digit decimal of an event
value; \texttt{p\_exact} is the declaration or the event's defining
relation with the primitives it uses. \texttt{regime} is the reserve's
position on the full payoff domain; \texttt{band\_status} says whether a
class reserve lies below, inside, or above each value band.
\texttt{low\_cost\_floor\_relation} and \texttt{ceiling\_relation} take
the values strict positive, strict negative, equality event, or
unresolved, from the fifty-digit evaluation. \texttt{candidate\_id} is
the attempt identity and \texttt{continuation\_id} the economic identity
defined above; \texttt{duplicate\_of} is empty or the representative's
candidate identity; \texttt{identity\_note} records
same-orders-different-atoms cases. \texttt{parameter\_set\_id} is the
canonical string of exact declarations (not a hash);
\texttt{institution\_id}, \texttt{information\_structure\_id}, and
\texttt{tie\_rule\_id} are the literal definitions of the sale rule, of
who observes what, and of the tie conventions. \texttt{order\_rule\_id}
is the canonical order profile; \texttt{price\_rule\_id} names the
pricing family, its pool sets, and the positive-entry price map.
\texttt{pricing\_family\_coverage} and
\texttt{exhaustive\_pricing\_search} are as declared in C.0.
\texttt{pool\_id}, \texttt{pool\_set\_definition},
\texttt{pool\_cutoff\_exact}, \texttt{price\_atom\_value},
\texttt{state\_H\_pool\_mass}, \texttt{state\_L\_pool\_mass},
\texttt{pool\_probability}, and \texttt{pool\_posterior} describe the
zero-entry atom when one exists and are \texttt{n/a} otherwise; a
positive-price plateau atom in a flat posterior tail is a different kind
of atom and is not written as a no-entry pool.
\texttt{preparation\_rule\_id} is the buyer's action by price region and
cost. \texttt{event\_id} and \texttt{event\_defining\_relation} identify
the node on the event grid. \texttt{preparation\_probability},
\texttt{admissible\_challenger\_probability},
\texttt{sale\_probability},
\texttt{two\_admissible\_bidders\_probability}, and
\texttt{high\_value\_ownership\_probability} are
\(\mathsf E,\mathsf A,\mathsf S,\mathsf C_2,\mathsf O_H\) of C.0;
\texttt{seller\_revenue} is \(\mathcal R_T\) and
\texttt{mean\_financial\_price} is \(\mathbb E P\), which must agree
with it. \texttt{pricing\_error}, \texttt{belief\_error}, and
\texttt{entry\_deviation\_gain} are the pointwise price, pooled-belief,
and buyer-optimality diagnostics;
\texttt{investor\_deviation\_gain\_estimate} is the maximal gain found
on the declared and refined order grids, and
\texttt{investor\_deviation\_gain\_upper\_bound} is a rigorous bound
when one exists (zero when an analytical bound covers the whole order
interval) and \texttt{n/a} otherwise, never zero by default.
\texttt{error\_budget} is the quadrature plus tail term.
\texttt{result\_status} uses the vocabulary of C.0;
\texttt{existence\_scope}, \texttt{uniqueness\_scope}, and
\texttt{search\_coverage\_scope} state, in words, what the row
establishes and what was searched. \texttt{rejection\_reason} carries
the breach witness of a rejected candidate; \texttt{unresolved\_reason}
the reason an open row is open. \texttt{run\_id} names the producing
exercise. In the range file, \texttt{search\_outcome} is one of the four
outcomes listed under Reporting, and
\texttt{global\_envelope\_certified} is false unless an additional
exhaustive argument establishes the full envelope.

Feed: Table 4 and the sale-design discussion. The research question
about seller-optimal discovery is not assigned an answer by this
exploratory file.

\subsection{C.6b. Price pooling at fixed orders}\label{oa-c-price-pools}

This exercise is the regression for A.11. It runs the declared cutoff
family at the price-pool reserve and the negative controls, and it is
the smallest test that the continuation representation above is right.

\textbf{Inputs.} The benchmark declaration with the reserve replaced by
the declared \texttt{pool\_reserve}, at strength \texttt{1.2}, full
orders \((1,-1)\), atomic costs, Laplace noise, and the cutoff grid
declared in C.0. The auction objects are computed in exact rational
arithmetic and by direct integration of \eqref{eq:oa-oa-50}, and both
must give \(t_0=t_L=g_L=0\), \(t_H=7\), \(g_H=3\).

\textbf{Method.} For each cutoff, build the complete pooled schedule
\eqref{eq:oa-oa-58} with its prescribed preparation policy and validate
it as a continuation. The pool masses \(\int_{-\infty}^ca_\theta\) are
computed in fifty-digit arithmetic from \eqref{eq:oa-oa-59} and by
segment-wise Gauss-Legendre integration, and the two must agree within
the independent-formula acceptance. The pooled posterior is
\eqref{eq:oa-oa-60} from the full preimage. Buyer optimality is checked
at the zero-price atom against the pooled posterior for both cost types,
and pointwise on the positive-price region on a flow mesh that starts at
the cutoff and includes every breakpoint; the row stores the minimum
slack in each case. Competitive pricing \eqref{eq:oa-oa-61} is checked
pointwise on a mesh spanning both regions and including the cutoff from
both sides, together with the residual identities \eqref{eq:oa-oa-62}
against direct conditional payoffs; the collision test requires a
positive gap between the positive-price range and the zero atom. The
upper plateau is recorded as a positive-entry atom with its own masses
and posterior. Investor deviations are scanned on the declared and
refined order grids for both signs and both types against the fixed
schedule. The global bound \eqref{eq:oa-oa-63} is evaluated in exact
rational arithmetic with an outward-rounded \(e^{2/b}\), and the row
stores the rational lower bounds on \(J_c\) and on \(U_\theta'\); the
numerical \(J_c\) from both types and the minimum of \(U_\theta'\) on a
two-hundred-point mesh are stored beside them as the implementation
check, and a numerical value below the rational bound is a breach. The
output quantities \eqref{eq:oa-oa-64} are computed in closed form and by
conditional integration; both are compared with each other and, at the
two endpoints, must agree to ten decimals with the values written into
the manifest as landmarks.

\textbf{Acceptance.} Every grid member passes every validation.
Preparation and revenue strictly decrease along the grid. The rational
bounds satisfy \(J_c\ge7/32\) and \(U'_\theta\ge143/1600\). The two
endpoints survive deduplication as distinct continuations, a repeated
construction of one endpoint from a second start merges into it with a
different \texttt{candidate\_id} and the same \texttt{continuation\_id},
and the orders-only key merges the two endpoints, which the ledger
records as the reason that key is forbidden. The controls \(c=-1\),
\(c=1\), and the raw-posterior buyer inside the \(c=0\) pool fail, and
each fails with the breach named in A.11. A control that does not fail
as required stops the exercise.

\textbf{Outputs.}

\begin{Shaded}
\begin{Highlighting}[]
\NormalTok{numerics/price\_pool\_regression.csv:}
\NormalTok{h, ell, p, rho, c\_L, c\_H, b, k, r, noise, cost, prior\_H, q\_H, q\_L,}
\NormalTok{t\_0, t\_H, t\_L, g\_H, g\_L, cutoff\_exact, cutoff\_decimal,}
\NormalTok{pool\_probability, state\_H\_pool\_mass, state\_L\_pool\_mass, pool\_posterior,}
\NormalTok{buyer\_slack\_zero\_price\_c\_L, buyer\_slack\_zero\_price\_c\_H,}
\NormalTok{buyer\_slack\_positive\_price\_c\_L\_min, buyer\_slack\_positive\_price\_c\_H\_min,}
\NormalTok{global\_bound\_J\_lower\_rational, global\_bound\_Uprime\_lower\_rational,}
\NormalTok{J\_numeric, min\_dU\_mesh,}
\NormalTok{preparation\_probability, sale\_probability, admissible\_challenger\_probability,}
\NormalTok{two\_admissible\_bidders\_probability, high\_value\_ownership\_probability,}
\NormalTok{seller\_revenue, mean\_financial\_price, closed\_form\_vs\_integration\_error,}
\NormalTok{pricing\_error, entry\_deviation\_gain, investor\_deviation\_gain\_estimate,}
\NormalTok{investor\_deviation\_gain\_upper\_bound, validation\_status, accepted,}
\NormalTok{expected\_failure\_reason, candidate\_id, continuation\_id, run\_id}
\end{Highlighting}
\end{Shaded}

\texttt{cutoff\_exact} is the symbolic cutoff and
\texttt{cutoff\_decimal} its thirty-two-digit expansion. Negative
controls and the duplicate start appear with \texttt{accepted\ =\ false}
and an \texttt{expected\_failure\_reason}; they are not accepted rows.
Feed: the registry keys \texttt{pool\_*} of C.8, Proposition A.10 in the
paper appendix, and A.11.

\subsection{C.6c. Exact reserve events and admissible-bid
outcomes}\label{oa-c-reserve-events}

This exercise runs the C.6 node solver, with the continuation identity
and event handling above, on the event grid alone. It leaves the full
sweep untouched and exists so that the event candidates and the outcome
measures can be checked without the hour-long scan.

\textbf{Inputs.} Both economies at both strengths, on the event grid of
C.6: \(0\), \(\ell\), \(r\), \(h\), the band edges in the class economy,
the declared reserve alternatives, \(p_L\) at every strength and \(p_H\)
where \(\ell<p_H<r\), each with the three one-sided offsets on both
sides, and the reserve \texttt{6.280} used as the sampled weak maximum
in the previous sweep.

\textbf{Method and acceptance.} At \(p_L\) in the weak economy, the row
must show the equality-event classification of the floor, a single
accepted full-order continuation whose status names the equality and the
tie rule rather than the strict inequality, \(\mathsf E=\rho\), sale
probability \(\rho/2\), \(\mathsf C_2=0\), and
\(\mathcal R_T=\rho p_L/2\) from both the row and a fifty-digit
evaluation; the three offsets below the event must classify the floor as
strictly positive and the three above as strictly negative. At \(p_H\)
in the strong economy, the row must show the equality-event
classification of the ceiling with a strictly positive floor and a
strictly positive full-order bound, \(\tau=M\), \(x^*=1\),
\(\alpha_H=1/2\), \(\alpha_L=e^{-2/b}/2\), and entry and revenue
agreeing with the fifty-digit evaluation; where \(p_H\) falls outside
its defining regime the node is an ordinary reserve node and must not be
classified as an event. At the sampled reserve \texttt{6.280},
\(\mathsf E=\rho\), sale \(\rho/2\), \(\mathsf C_2=0\), and
\(\mathcal R_T=\rho\,(6.280)/2\): preparation there does not mean two
admissible bidders, because the incumbent is excluded. On every row the
union identity for \(\mathsf S\) and the ordering
\(0\le\mathsf C_2\le\mathsf A\le\mathsf E\le1\) hold, and the payoff
oracle error is within tolerance. Deduplication follows the C.6 rule,
and every count in the manifest comes from the ledger.

\textbf{Outputs.} \texttt{numerics/reserve\_events.csv}, with the
continuation schema of C.6 and an added \texttt{event\_classification}
text column, and \texttt{numerics/reserve\_event\_ranges.csv} with the
range schema of C.6. Feed: the exact-event paragraphs of the paper
appendix and the online summary of C.6. Event candidates improve sampled
maxima; they prove neither global optimality nor envelope completeness,
and a node without an accepted candidate is a node where the declared
searches found none.

\subsection{C.7. Bargaining weights and the payment
property}\label{oa-c-bargaining}

\textbf{Inputs.} Use the benchmark \(h,\ell\) and weak/strong incumbent
distributions, but set the bargaining institution's reserve to zero as
specified in Proposition A.8. Evaluate seller weights from \texttt{0} to
\texttt{0.99} by \texttt{0.01}, including the exact midpoint
\texttt{0.5}. The endpoint \(\eta=1\) can be stored separately as a
payment-stage limit, not as a positive-cost entry equilibrium.

\textbf{Objects and verification.} Compute \(T_\eta(R,v)\) and the
winning challenger's profit pointwise, then integrate over the
incumbent. Independently compare the results with \eqref{eq:oa-oa-44},
including no-entry revenue \(\eta r/2\). Verify feasibility of each
transfer between fallback and winning value and the Nash first-order
condition for interior weights. Compute differences between strong and
weak distributions for \(\Delta_\eta\) and \(G_{\theta,\eta}\). The
spread difference must change sign at the midpoint, with zero at that
exact weight. This is an acquisition-stage exercise; do not append a
trader or entry prediction without solving the new continuation game.

\textbf{Outputs.}

\begin{Shaded}
\begin{Highlighting}[]
\NormalTok{figures\_data/bargaining.csv:}
\NormalTok{eta, r, t\_0\_eta, t\_H\_eta, t\_L\_eta, Delta\_eta,}
\NormalTok{G\_H\_eta, G\_L\_eta, spread\_strength\_difference,}
\NormalTok{profit\_H\_strength\_difference, profit\_L\_strength\_difference,}
\NormalTok{transfer\_feasibility\_error, integration\_error, status}
\end{Highlighting}
\end{Shaded}

Feed: Figure 4 and Proposition A.8. Parameter values and the change of
institution are stated in the figure caption.

\subsection{C.8. Quantity definitions and
provenance}\label{oa-c-registry}

Every quantitative placeholder in the two manuscripts resolves to one
row of the \href{../numerics/quantity_registry.csv}{quantity registry},
generated from the \href{quantity_manifest.csv}{quantity manifest} and
the validated exercise outputs in C.1 to C.7. The manifest specifies
each key's definition, producing exercise, source file, row selector,
display rule, units, and exact value for a declared input. The registry
records the resulting value, outward bounds for an enclosure, evidence
status, and the source row used. The complete key list and definitions
are in the distributed
\href{../replication/quantity_dictionary.md}{quantity dictionary},
generated from the same manifest. This subsection states their common
rules.

\textbf{Rules.} A declared input is echoed with status \texttt{input}
and printed as its exact decimal. A derived quantity is filled only from
an accepted source row selected by the full parameter declaration,
information regime, cost and noise law, continuation identity, and
branch label, never by a rounded scalar. The selector must match exactly
one row. A difference records both rows used to form it. The row's
result status must permit the statement the manuscript makes with it: an
analytical claim needs an analytical row, a certified interval needs
every certificate predicate, and a numerical diagnostic is printed as
one. An unresolved source leaves the placeholder unfilled with its
reason, and the substitution step refuses to produce a filled manuscript
in that state. Mathematical constants, equation and result labels,
bibliographic identifiers, and schema labels are literal text, not
registry quantities. No renderer replaces a registry value with a cached
or typed number.

\textbf{Units and rounding.} Probabilities are stored as probabilities
and displayed by \texttt{decimal\_6}; \texttt{percent\_integer} converts
a declared accuracy to a percentage at display time only. Margins and
bounds use \texttt{scientific\_10}; thresholds and standardized
distances use \texttt{decimal\_9}. Enclosures use
\texttt{outward\_interval\_8} or \texttt{outward\_interval\_10},
rounding the lower endpoint down and the upper endpoint up; one-sided
bounds use \texttt{lower\_bound\_10}, \texttt{lower\_bound\_12}, or
\texttt{upper\_bound\_12}, rounding conservatively. Rounding is applied
after every validation on the unrounded value. A displayed lower bound
that supports a strict inequality must remain strictly positive after
rounding; if outward rounding would print a zero, more digits are
printed rather than rounding inward. Model units are the per-share
normalization of A.1.

\textbf{Quantity families.}

{\def\LTcaptype{none} % do not increment counter
\begin{longtable}[]{@{}
  >{\raggedright\arraybackslash}p{(\linewidth - 4\tabcolsep) * \real{0.3333}}
  >{\raggedright\arraybackslash}p{(\linewidth - 4\tabcolsep) * \real{0.3333}}
  >{\raggedright\arraybackslash}p{(\linewidth - 4\tabcolsep) * \real{0.3333}}@{}}
\toprule\noalign{}
\begin{minipage}[b]{\linewidth}\raggedright
Quantity group
\end{minipage} & \begin{minipage}[b]{\linewidth}\raggedright
Exercise
\end{minipage} & \begin{minipage}[b]{\linewidth}\raggedright
Units
\end{minipage} \\
\midrule\noalign{}
\endhead
\bottomrule\noalign{}
\endlastfoot
Declared inputs & C.0 & exact input \\
Auction payoffs & C.1 & model units \\
Equilibrium and control outcomes & C.1 & probability; model units \\
Welfare and revenue differences & C.1 & model units \\
Sufficient-condition margins & C.1, C.3, C.4 & model units \\
Thresholds and posterior bounds & C.2 & model units; probability \\
Distributional and moderate examples & C.1, C.4, C.5 & probability;
normalized distance \\
Complementary signals & C.3 & probability; percentage \\
Certified asymmetric nodes & C.2 & order; probability; model units \\
Class-economy reserve comparison & C.6 & probability; model units \\
Price-pool endpoints & C.6b & probability; model units \\
\end{longtable}
}

The basic definitions the families share are
\(\mathsf E=(\bar e_H+\bar e_L)/2\), \(\mathsf O_H=\bar e_H/2\) on the
benchmark support and the allocation-event value elsewhere,
\(\mathcal R_T=t_0+[\bar e_H(t_H-t_0)+\bar e_L(t_L-t_0)]/2\), the
margins \eqref{eq:oa-oa-77} and \eqref{eq:oa-oa-29}, the thresholds of
A.5, the certificate predicates of B.6, and the outcome measures of C.0.
Within a parameter set, \texttt{r\_weak}, \texttt{r\_strong}, and
\texttt{r\_collapse} select the declared strengths; the certified keys
refer to the three declared nodes in order, not to roots chosen from a
larger search; the reserve keys select on \texttt{value\_law},
\texttt{epsilon\_V}, \texttt{r}, and \texttt{p} exactly; the price-pool
keys select on the symbolic cutoff. The minimum-margin keys are a
convenience for acceptance and do not replace the recorded components. A
figure is never the source of a scalar, and an apparent curve is never
used to fill a missing certificate.

\textbf{Production order.} The benchmark, moderate-value, and
complementary-signal inputs and margins come first, then the threshold
boundaries, the benchmark entry, control, ownership, and welfare
quantities, and the certified brackets; these determine which theorem
regions the text can populate. The correspondence and reserve sweeps,
the price-pool regression, and the event exercise follow.
\texttt{numerics/check\_registry.py} regenerates the registry in fresh
processes under several ambient decimal precisions and import orders and
requires identical output, so the printed values do not depend on the
environment that produced them.

\section{D. Institutional pilot and empirical design}\label{oa-d}

The proposed pilot identifies publicly visible sale opportunities in
which a prospective buyer can still decide whether to investigate while
target shares trade. This appendix specifies eligibility, event coding,
and the evidence needed to establish the decision interval. It reports a
research design, without an extracted sample or estimated effects.

\subsection{D.1. The decision interval and sample-selection
rule}\label{oa-d-selection}

The pilot establishes the institution required by the model before
attempting to estimate its causal mechanism. The unit of observation is
a prospective buyer's participation decision within a publicly visible
sale process for a listed target. A process is eligible when
contemporaneous public information identifies a plausible acquisition
opportunity, the target continues to trade, and at least one
economically meaningful participation decision by a potential challenger
remains open. Public visibility and the buyer's later decision must be
separately dated.

A public strategic review, disclosed approach, or open bidding contest
is a candidate starting point, not automatic inclusion. The public
record must show that a prospective buyer could still decide whether to
investigate or submit a substantive proposal. A transaction announced
only after the buyer set was fixed does not qualify merely because the
target traded during its earlier confidential negotiations. Likewise, a
named firm mentioned by an adviser is a potential buyer, not an entrant;
an NDA alone need not establish costly acquisition preparation. I retain
each of these events but distinguish their economic content.

Selection proceeds without conditioning on whether a later bidder wins,
whether a target return has a particular sign, or whether the process
appears to support the theory. Construct a process-level screening log
from the eligible filing and announcement universe, record inclusion and
exclusion decisions and their evidence, and retain uncertain cases for
adjudication. Multiple processes for the same target receive separate
process identifiers when the record documents a termination and restart.
No cases are selected and no sample count or estimate is reported in
this design.

\subsection{D.2. Information available at the time, not only at the
filing date}\label{oa-d-information}

I distinguish the date an event occurred from the date its content first
became public. A definitive proxy can describe a private meeting
retrospectively. That description supports the occurrence of the
meeting, but does not put it in the stock market's information set on
the meeting date. For every economically relevant event, record the
occurrence date or bounded date range, the first verified disclosure
date, the public source, and the specificity of the disclosed
information.

The candidate modeled interval begins when the opportunity and relevant
competitive threat are publicly interpretable and ends at the
challenger's material participation decision. If the rival's identity,
financing position, or proposed terms become public later, those facts
enter the information set only then. An unverified rumor is coded as a
rumor with its source and uncertainty, not as an established public bid.
If the chronology provides only a month or a relative ordering, preserve
that precision rather than inventing a daily date.

A process can contain more than one candidate interval. The pilot
retains the chronology and identifies each interval explicitly; it does
not select the interval producing the strongest return--entry
association. Confidential diligence by the challenger may follow public
visibility even when its identity remains private to the market. In that
case the record must support an interpretation of what match-related
information a public investor could hold, rather than presuming that the
market knew the eventual entrant.

\subsection{D.3. Sources and event coding}\label{oa-d-coding}

The primary chronology is the target's merger-background discussion in
definitive proxy or tender-offer filings, supplemented by current-report
filings, transaction exhibits, contemporaneous company releases, and
time-stamped public reporting. SEC accession numbers and exact source
locations accompany every coded factual claim. A filing's narrative is
evidence, not an instruction to treat every party as having the same
information. Contradictory dates or actor descriptions remain visible
until resolved.

I use the following event categories in the pilot. An initial approach
identifies who initiated contact and whether it was solicited. Public
visibility identifies the first disclosure of the acquisition
opportunity and the competitive facts revealed. Buyer contact records
outreach and replies without equating them with entry. Confidentiality
and data access record when information became available to the buyer.
Diligence entry records a substantiated commitment to evaluation,
including its stated scope. A first substantive proposal records price,
consideration, financing, and diligence conditions. A revised proposal
records what changed and which new information preceded it. Withdrawal
records the actor making the decision and the reason when stated. Final
selection records the board's choice, agreement, or process termination.
Announced and actual deadlines are separate events.

The event ledger has a row for each event--actor link. The following
schema is a data specification, not an extracted dataset:

\begin{Shaded}
\begin{Highlighting}[]
\NormalTok{pilot/processes.csv:}
\NormalTok{process\_id, target\_id, target\_name, listing\_status,}
\NormalTok{process\_start\_date\_lower, process\_start\_date\_upper,}
\NormalTok{public\_visibility\_date\_lower, public\_visibility\_date\_upper,}
\NormalTok{process\_end\_date\_lower, process\_end\_date\_upper,}
\NormalTok{lead\_buyer\_id, process\_status, eligible, exclusion\_reason,}
\NormalTok{information\_interval\_established, unresolved\_issue}

\NormalTok{pilot/events.csv:}
\NormalTok{process\_id, event\_id, actor\_id, actor\_role, event\_type,}
\NormalTok{event\_date\_lower, event\_date\_upper, event\_date\_precision,}
\NormalTok{first\_public\_date\_lower, first\_public\_date\_upper,}
\NormalTok{public\_information\_content, contemporaneous\_or\_retrospective,}
\NormalTok{price\_or\_range, consideration, financing\_condition,}
\NormalTok{diligence\_condition, binding\_status, initiation\_actor,}
\NormalTok{withdrawal\_decision\_actor, reported\_reason,}
\NormalTok{accession\_or\_source\_id, source\_locator, verbatim\_evidence,}
\NormalTok{source\_publication\_time, evidence\_confidence, adjudication\_status}

\NormalTok{pilot/intervals.csv:}
\NormalTok{process\_id, interval\_id, challenger\_id,}
\NormalTok{interval\_start\_lower, interval\_start\_upper,}
\NormalTok{participation\_decision\_lower, participation\_decision\_upper,}
\NormalTok{public\_opportunity\_evidence, public\_incumbent\_evidence,}
\NormalTok{participation\_evidence, information\_complementarity\_interpretation,}
\NormalTok{ordering\_verified, eligibility\_decision, unresolved\_issue}

\NormalTok{pilot/adjudication.csv:}
\NormalTok{record\_id, field, initial\_value, alternative\_value,}
\NormalTok{conflict\_type, supporting\_sources, resolution, resolver,}
\NormalTok{resolution\_date, unresolved\_reason}
\end{Highlighting}
\end{Shaded}

An absent field means missing, not applicable, or not stated, with a
corresponding reason; it is not a zero price, an absence of a buyer, or
a voluntary withdrawal. For price ranges, store both endpoints rather
than the midpoint alone. A proposal's classification as formal or
informal follows the source's description and the coded contractual
features; no uniform price threshold or arbitrary document label is used
to manufacture a distinction across processes. Quotations are checked
against the source location, and summaries are stored separately from
the evidence text.

\subsection{D.4. What would establish institutional
relevance}\label{oa-d-relevance}

The pilot supports the timing if it documents a public opportunity,
trading during that opportunity, and a subsequent participation decision
that remained open. It supports the preparation margin if a buyer had to
commit resources or obtain additional target information before
submitting an executable bid. It supports an incremental-information
interpretation if the relevant uncertainty was not already fully
resolved by the buyer's public statements and the record leaves room for
distinct target-side and buyer-side knowledge.

None of these observations alone establishes that prices caused entry. A
target price rising before a competing bid can anticipate that bid. A
private process described after announcement does not show that the
public market knew it beforehand. A buyer's later success does not prove
that it entered after receiving a favorable price signal. Such cases may
illuminate institutions while remaining outside the mechanism's direct
empirical evidence.

The first output is a sourced chronology and an accounting of eligible
decision intervals, with exclusions and ambiguous cases retained. The
next output is a short institutional section describing which elements
of the modeled game occur in the verified intervals. The existing
process illustration in the manuscript is not a substitute for this
selection and coding exercise.

\subsection{D.5. Measurement and eventual empirical
tests}\label{oa-d-measurement}

The participation outcome is commitment to investigation or substantive
bidding, measured at the most precise date the evidence supports. Public
bid counts are a separate, narrower outcome. High-quality ownership in
the theory is a structural value concept; a winning challenger is
observable, but winning alone is not evidence that its underlying
acquisition value was high. Any empirical proxy for match quality must
be separately justified.

Incumbent strength is measured using information available before the
challenger's decision. A final offer reflects subsequent entry and
cannot serve as a predetermined initial threat. Candidate descriptive
measures include an initial disclosed proposal, pre-existing financing
capacity, and publicly measurable buyer--target complementarity. Their
interpretation must respect whether the measure changes competitive
strength, common acquisition value, financing constraints, or several
objects at once.

Returns, trading activity, and other information measures are
timestamped before the outcome. Define the public-information cutoff
before constructing financial windows, and retain alternative date
bounds when the event time is interval-censored. A stock return is not
itself a measure of how much information was revealed. The empirical
design must distinguish a change in the expected level of sale proceeds
from a change in the market's information about challenger quality.

A targeted causal design would shift information transmission without
independently changing acquisition values, preparation costs, or
financing. No such instrument is asserted here. Alternatively, a
structural exercise could jointly model anticipated entry and learning
from prices, using the final institution-specific theorem to identify
discriminating restrictions. The pilot determines which route is
institutionally credible before either exercise is undertaken.

\section{E. Reproducibility and the numerical boundary}\label{oa-e}

The replication package separates independent node verification, the
numerical exercises, and presentation. This section describes their
inputs, recorded environments, and reproducibility checks.

\subsection{E.1. Distributed verification and execution
environment}\label{oa-e-environment}

The verification directory contains independent node calculations,
exploratory asymmetric searches, and interval certificates. The
documented execution environment is Python 3.13.5, NumPy 2.3.5, SciPy
1.17.0, and mpmath 1.3.0. These are metadata for the distributed
results, not a claim that every other compatible environment fails. Each
new execution records its actual interpreter, package versions,
operating system, arithmetic precision, parameter declarations,
tolerances, and source-file hashes.

The numerical layer under \texttt{numerics/} implements the exercises in
\hyperref[oa-c]{Appendix C}, including the searches over incumbent
strengths and reserves, with the coverage limits stated there. Accepted
output rows supply the quantity registry, tables, and figures, and each
exercise writes a run manifest under \texttt{numerics/manifests/} with
its inputs, method, tolerances, software versions, the SHA-256 of every
output, and its pass or fail result. The verification gate refuses any
output whose hash no longer matches its manifest, so a rerun of the
producing exercise is the only way to change a number. The empirical
pilot in \hyperref[oa-d]{Appendix D} remains a coding specification and
produces nothing in this build.

Each displayed exercise has one producer, run from the repository root
with the project interpreter. The output schemas in Appendix C identify
its files, including both C.2 attempt ledgers. Commands are:

\begin{Shaded}
\begin{Highlighting}[]
\BuiltInTok{source}\NormalTok{ .venv/bin/activate}
\BuiltInTok{export} \VariableTok{OPENBLAS\_NUM\_THREADS}\OperatorTok{=}\NormalTok{1 }\VariableTok{OMP\_NUM\_THREADS}\OperatorTok{=}\NormalTok{1 }\VariableTok{MKL\_NUM\_THREADS}\OperatorTok{=}\NormalTok{1}
\BuiltInTok{export} \VariableTok{NUMEXPR\_NUM\_THREADS}\OperatorTok{=}\NormalTok{1}
\VariableTok{PY}\OperatorTok{=}\NormalTok{.venv/bin/python}
\VariableTok{$PY}\NormalTok{ numerics/exercises/c1\_baseline.py}
\VariableTok{$PY}\NormalTok{ numerics/exercises/c2\_correspondence.py }\AttributeTok{{-}{-}workers}\OperatorTok{=}\NormalTok{20}
\VariableTok{$PY}\NormalTok{ numerics/exercises/c3\_signals.py}
\VariableTok{$PY}\NormalTok{ numerics/exercises/c4\_moderate.py}
\VariableTok{$PY}\NormalTok{ numerics/exercises/c5\_noise.py}
\VariableTok{$PY}\NormalTok{ numerics/exercises/c6\_reserve.py }\AttributeTok{{-}{-}workers}\OperatorTok{=}\NormalTok{20}
\VariableTok{$PY}\NormalTok{ numerics/exercises/c6b\_price\_pools.py}
\VariableTok{$PY}\NormalTok{ numerics/exercises/c6c\_reserve\_events.py }\AttributeTok{{-}{-}workers}\OperatorTok{=}\NormalTok{8}
\VariableTok{$PY}\NormalTok{ numerics/exercises/c7\_bargaining.py}
\end{Highlighting}
\end{Shaded}

The broad exercises use twenty workers on the replication VM, and the
event exercise uses eight. The thread settings prevent numerical
libraries from multiplying process concurrency. A smaller worker count
changes runtime without changing the scientific grid or tolerances.
Reduced diagnostic runs use \texttt{-\/-quick} and a separate
\texttt{-\/-out} directory where supported. They do not replace the
declared release corpus. The presentation layer then runs in this order:

\begin{Shaded}
\begin{Highlighting}[]
\VariableTok{$PY}\NormalTok{ numerics/check\_registry.py}
\VariableTok{$PY}\NormalTok{ numerics/registry.py}
\VariableTok{$PY}\NormalTok{ numerics/substitute.py}
\VariableTok{$PY}\NormalTok{ numerics/render/render\_all.py}
\VariableTok{$PY}\NormalTok{ numerics/render/latex.py}
\VariableTok{$PY}\NormalTok{ numerics/verify.py }\AttributeTok{{-}{-}final}
\end{Highlighting}
\end{Shaded}

Two entry points wrap these sequences. \texttt{make\ peer-release}
validates the existing completed calculation outputs, runs the tests,
regenerates presentation, typesets, runs the final gate, and assembles
the deliverables after checking a matching visual-inspection record. For
a new revision, the release script's \texttt{-\/-build-only} option
produces the files for inspection and \texttt{-\/-package-only}
validates and packages them afterward. \texttt{make\ peer-reproduce}
runs the same declared sequence in a fresh working and output directory,
without any accepted result carried over from a previous run, and
compares the canonical numerical outputs, exact input declarations,
certificate validity, resolved references, and extracted manuscript text
with the release. The release runner fixes creation-date metadata with
\texttt{SOURCE\_DATE\_EPOCH}; the comparison requires identical
substantive content rather than relying on byte-identical PDFs.
\texttt{replication/README.md} lists every producer command and the role
of every file.

Registry calculations use a local decimal context with precision 60, so
their serialized values do not depend on ambient arithmetic precision or
import order; the registry check runs the generator in fresh processes
at several ambient precisions and import orders and requires identical
output. The final verification includes the registry manifest, numerical
acceptance checks, and strict substitution. Source manuscripts remain
editable; filled Markdown, LaTeX, tables, and figures are regenerated.

The \texttt{verification/} directory is the independent seed delivered
before the numerical layer existed. The numerical exercises do not
import its calculations. The independent certificate audit compares its
separately rerun results with the local interval port. To rerun its
checks in isolation, use:

\begin{Shaded}
\begin{Highlighting}[]
\ExtensionTok{python} \AttributeTok{{-}m}\NormalTok{ pip install }\AttributeTok{{-}r}\NormalTok{ requirements.txt}
\ExtensionTok{python}\NormalTok{ review\_checks.py }\AttributeTok{{-}{-}out}\NormalTok{ results }\AttributeTok{{-}{-}mode}\NormalTok{ core}
\ExtensionTok{python}\NormalTok{ review\_checks.py }\AttributeTok{{-}{-}out}\NormalTok{ results }\AttributeTok{{-}{-}mode}\NormalTok{ signals}
\ExtensionTok{python}\NormalTok{ explore\_asymmetric.py}
\end{Highlighting}
\end{Shaded}

Run these commands from an isolated copy of \texttt{verification/} and
retain the distributed result files separately from newly generated
results. For the certificate, run the following from the repository
root:

\begin{Shaded}
\begin{Highlighting}[]
\FunctionTok{bash}\NormalTok{ audit/peer\_polish/reference\_seed/run\_seed.sh}
\end{Highlighting}
\end{Shaded}

This guarded runner executes a byte-identical seed copy and refuses
optimized Python or \texttt{PYTHONOPTIMIZE}, because the seed uses
assertions for acceptance. Record a failure as a failure; do not change
a tolerance or catch an exception solely to obtain an accepted output.
The interval certificate requires outward interval arithmetic and its
analytical cover. Replacing it with ordinary floating-point quadrature
changes the evidentiary status.

\subsection{E.2. Contents and scope of the distributed
files}\label{oa-e-files}

The core verifier computes auction expectations, theorem margins,
full-order objects, threshold distinctions, moderate values, noise-law
comparisons, and the specified class-value reserve alternatives. Its
signal mode uses the conditional laws of Appendix A.7. The exploratory
asymmetric search identifies candidate roots and checks numerical
deviations. The certificate evaluator uses the antiderivatives and
derivative cover of \hyperref[oa-b]{Appendix B} to establish selected
continuous-action equilibria.

The distributed directory contains:

\begin{Shaded}
\begin{Highlighting}[]
\NormalTok{verification/}
\NormalTok{  review\_checks.py}
\NormalTok{  explore\_asymmetric.py}
\NormalTok{  certify\_asymmetric.py}
\NormalTok{  requirements.txt}
\NormalTok{  results/}
\NormalTok{    core\_results.json}
\NormalTok{    deviations.csv}
\NormalTok{    two\_signal\_results.json}
\NormalTok{    two\_signal\_deviations.csv}
\NormalTok{    asymmetric\_candidates.json}
\NormalTok{    asymmetric\_interval\_certificates.json}
\NormalTok{    certificate\_run\_log.txt}
\end{Highlighting}
\end{Shaded}

The core results hold primitive vectors, inequality margins, threshold
calculations, candidate outcomes, and reserve comparisons. The deviation
ledger records tested state--order alternatives and their payoffs. The
signal results hold private-information margins, joint-probability
checks, and the associated deviation ledger. The asymmetric candidate
file is exploratory: its root and local-search output is not an
exhaustive equilibrium enumeration. The interval file contains the
root-sign enclosures, order-cover bounds, and entry intervals that
complete the selected computer-assisted existence arguments. The
certificate log provides the readable arithmetic output; the complete
interval endpoints, not their displayed midpoints, are the numerical
proof objects.

These files support the statements at their documented parameter points.
They do not supply an optimizer over sale terms, a first-price bidding
equilibrium, or the full set of mixed trading equilibria. The exercises
in \hyperref[oa-c]{Appendix C} extend numerical coverage while
preserving those limits. Their results are reported separately from the
independent node checks.

\subsection{E.3. Implementation contract for the full
exercises}\label{oa-e-contract}

The numerical layer is written as pure functions with explicit parameter
arguments and immutable returned records. Its modules keep the layers
apart: \texttt{auction.py} (payoffs on the full reserve domain, with the
direct integration oracle), \texttt{noise.py} and
\texttt{information.py} (densities, posteriors, price schedules),
\texttt{deviations.py} and \texttt{quadrature.py} (unilateral-deviation
payoffs with analytic Laplace tails), \texttt{search.py} (candidate
search), \texttt{validation.py} and \texttt{continuations.py}
(independent validation, price atoms, continuation identity, and the
wide schema), \texttt{reserve\_events.py} (exact events and regime
classification), \texttt{certificates.py} (the interval port of
\hyperref[oa-b]{Appendix B}), \texttt{thresholds.py},
\texttt{signals.py}, and \texttt{render/} (figures and tables).
\texttt{params.py} holds the exact decimal declarations of C.0 as
strings and the \texttt{Controls} record of tolerances; exercises parse
them as decimals and never through a binary float where a certificate
depends on them. A search routine returns candidates and unresolved
nodes; only the validation layer assigns \texttt{accepted}. No renderer
solves an equilibrium, changes a parameter, or drops a branch.

The only boundary to the presentation layer is the validated CSV output
specified in \hyperref[oa-c]{Appendix C}, written through one writer
that prints floats by their shortest round-trip representation and marks
a missing value \texttt{n/a}. Tables are joined on the complete
parameter declaration and branch or continuation label, never on a
rounded strength or price. The extensions table is assembled from the
distributional, moderate-value, and signal outputs, with fields that do
not apply left as not applicable; a missing signal bound or an
uncomputed welfare statistic is never filled with zero. The reserve
table uses the declared comparisons; the found-continuation ranges and
the event candidates are reported online under the heading fixed in C.6.

The scalar registry is generated from validated rows and exact input
declarations. The substitution stage checks that every placeholder key
is defined once, that its source row exists and is unique under the
declared selector, and that its result status permits the intended
statement. It fails before creating a filled manuscript if a required
value remains open. Certificates are substituted as outward intervals
and conservative lower bounds. Ordinary diagnostics carry their
specified rounding without being promoted to certified enclosures.

Every exercise writes a manifest with inputs, method, tolerances,
software versions, output hashes, and pass or fail outcomes;
\texttt{numerics/verify.py} reads every manifest and fails on a missing
pass, a changed hash, a failed revalidation of the declared nodes, a
failed certificate, or an open placeholder. Failed searches, negative
controls, duplicate starts, and arithmetic retries remain in the record
with their reasons. A reader can therefore distinguish a rejected
candidate from a numerical failure, an unresolved search from a
nonexistence argument, and an accepted node from an exhaustive
characterization.

\subsection{E.4. Figure and table rendering}\label{oa-e-rendering}

Figures have no internal titles. Captions state the economic comparison,
parameter specification, information regime, units, and result-status
interpretation. Multi-panel figures use the panel labels specified in
the manuscript. The correspondence figure shows separate branches and
breaks lines at unresolved nodes. Certified nodes retain interval bars
even where the intervals are visually small. A uniqueness boundary is
labeled by the theorem it satisfies, not by the first root found by a
solver.

Tables preserve the difference between equilibrium comparisons and
fixed-profile controls. A frozen-profile row can contain profitable
investor deviations without contradicting its purpose; it is not labeled
an equilibrium. Welfare columns distinguish target revenue from
allocation surplus net of preparation costs. Reserve ranges are
described as ranges across continuations found unless an exhaustive
argument establishes more. No selected branch, matched mean price, or
apparent hump replaces the underlying equilibrium checks.

\protect\phantomsection\label{refs}
\begin{CSLReferences}{1}{1}
\bibitem[\citeproctext]{ref-DowGoldsteinGuembel2017}
Dow, James, Itay Goldstein, and Alexander Guembel. 2017. {``Incentives
for Information Production in Markets Where Prices Affect Real
Investment.''} \emph{Journal of the European Economic Association} 15
(4): 877--909. \url{https://doi.org/10.1093/jeea/jvw023}.

\end{CSLReferences}

\end{document}
